\documentclass[
  doc,
  floatsintext,
  longtable,
  a4paper,
  nolmodern,
  notxfonts,
  notimes,
  12pt,
  colorlinks=true,linkcolor=blue,citecolor=blue,urlcolor=blue]{apa7}

\usepackage{amsmath}
\usepackage{amssymb}

\geometry{inner=1in, outer=1in}
\fancyhfoffset[LE,RO]{0cm}


\usepackage[bidi=default]{babel}
\babelprovide[main,import]{spanish}


% get rid of language-specific shorthands (see #6817):
\let\LanguageShortHands\languageshorthands
\def\languageshorthands#1{}

\RequirePackage{longtable}
\RequirePackage{threeparttablex}

\makeatletter
\renewcommand{\paragraph}{\@startsection{paragraph}{4}{\parindent}%
	{0\baselineskip \@plus 0.2ex \@minus 0.2ex}%
	{-.5em}%
	{\normalfont\normalsize\bfseries\typesectitle}}

\renewcommand{\subparagraph}[1]{\@startsection{subparagraph}{5}{0.5em}%
	{0\baselineskip \@plus 0.2ex \@minus 0.2ex}%
	{-\z@\relax}%
	{\normalfont\normalsize\bfseries\itshape\hspace{\parindent}{#1}\textit{\addperi}}{\relax}}
\makeatother




\usepackage{longtable, booktabs, multirow, multicol, colortbl, hhline, caption, array, float, xpatch}
\usepackage{subcaption}


\renewcommand\thesubfigure{\Alph{subfigure}}
\setcounter{topnumber}{2}
\setcounter{bottomnumber}{2}
\setcounter{totalnumber}{4}
\renewcommand{\topfraction}{0.85}
\renewcommand{\bottomfraction}{0.85}
\renewcommand{\textfraction}{0.15}
\renewcommand{\floatpagefraction}{0.7}

\usepackage{tcolorbox}
\tcbuselibrary{listings,theorems, breakable, skins}
\usepackage{fontawesome5}

\definecolor{quarto-callout-color}{HTML}{909090}
\definecolor{quarto-callout-note-color}{HTML}{0758E5}
\definecolor{quarto-callout-important-color}{HTML}{CC1914}
\definecolor{quarto-callout-warning-color}{HTML}{EB9113}
\definecolor{quarto-callout-tip-color}{HTML}{00A047}
\definecolor{quarto-callout-caution-color}{HTML}{FC5300}
\definecolor{quarto-callout-color-frame}{HTML}{ACACAC}
\definecolor{quarto-callout-note-color-frame}{HTML}{4582EC}
\definecolor{quarto-callout-important-color-frame}{HTML}{D9534F}
\definecolor{quarto-callout-warning-color-frame}{HTML}{F0AD4E}
\definecolor{quarto-callout-tip-color-frame}{HTML}{02B875}
\definecolor{quarto-callout-caution-color-frame}{HTML}{FD7E14}

%\newlength\Oldarrayrulewidth
%\newlength\Oldtabcolsep


\usepackage{hyperref}



\usepackage{color}
\usepackage{fancyvrb}
\newcommand{\VerbBar}{|}
\newcommand{\VERB}{\Verb[commandchars=\\\{\}]}
\DefineVerbatimEnvironment{Highlighting}{Verbatim}{commandchars=\\\{\}}
% Add ',fontsize=\small' for more characters per line
\usepackage{framed}
\definecolor{shadecolor}{RGB}{241,243,245}
\newenvironment{Shaded}{\begin{snugshade}}{\end{snugshade}}
\newcommand{\AlertTok}[1]{\textcolor[rgb]{0.68,0.00,0.00}{#1}}
\newcommand{\AnnotationTok}[1]{\textcolor[rgb]{0.37,0.37,0.37}{#1}}
\newcommand{\AttributeTok}[1]{\textcolor[rgb]{0.40,0.45,0.13}{#1}}
\newcommand{\BaseNTok}[1]{\textcolor[rgb]{0.68,0.00,0.00}{#1}}
\newcommand{\BuiltInTok}[1]{\textcolor[rgb]{0.00,0.23,0.31}{#1}}
\newcommand{\CharTok}[1]{\textcolor[rgb]{0.13,0.47,0.30}{#1}}
\newcommand{\CommentTok}[1]{\textcolor[rgb]{0.37,0.37,0.37}{#1}}
\newcommand{\CommentVarTok}[1]{\textcolor[rgb]{0.37,0.37,0.37}{\textit{#1}}}
\newcommand{\ConstantTok}[1]{\textcolor[rgb]{0.56,0.35,0.01}{#1}}
\newcommand{\ControlFlowTok}[1]{\textcolor[rgb]{0.00,0.23,0.31}{\textbf{#1}}}
\newcommand{\DataTypeTok}[1]{\textcolor[rgb]{0.68,0.00,0.00}{#1}}
\newcommand{\DecValTok}[1]{\textcolor[rgb]{0.68,0.00,0.00}{#1}}
\newcommand{\DocumentationTok}[1]{\textcolor[rgb]{0.37,0.37,0.37}{\textit{#1}}}
\newcommand{\ErrorTok}[1]{\textcolor[rgb]{0.68,0.00,0.00}{#1}}
\newcommand{\ExtensionTok}[1]{\textcolor[rgb]{0.00,0.23,0.31}{#1}}
\newcommand{\FloatTok}[1]{\textcolor[rgb]{0.68,0.00,0.00}{#1}}
\newcommand{\FunctionTok}[1]{\textcolor[rgb]{0.28,0.35,0.67}{#1}}
\newcommand{\ImportTok}[1]{\textcolor[rgb]{0.00,0.46,0.62}{#1}}
\newcommand{\InformationTok}[1]{\textcolor[rgb]{0.37,0.37,0.37}{#1}}
\newcommand{\KeywordTok}[1]{\textcolor[rgb]{0.00,0.23,0.31}{\textbf{#1}}}
\newcommand{\NormalTok}[1]{\textcolor[rgb]{0.00,0.23,0.31}{#1}}
\newcommand{\OperatorTok}[1]{\textcolor[rgb]{0.37,0.37,0.37}{#1}}
\newcommand{\OtherTok}[1]{\textcolor[rgb]{0.00,0.23,0.31}{#1}}
\newcommand{\PreprocessorTok}[1]{\textcolor[rgb]{0.68,0.00,0.00}{#1}}
\newcommand{\RegionMarkerTok}[1]{\textcolor[rgb]{0.00,0.23,0.31}{#1}}
\newcommand{\SpecialCharTok}[1]{\textcolor[rgb]{0.37,0.37,0.37}{#1}}
\newcommand{\SpecialStringTok}[1]{\textcolor[rgb]{0.13,0.47,0.30}{#1}}
\newcommand{\StringTok}[1]{\textcolor[rgb]{0.13,0.47,0.30}{#1}}
\newcommand{\VariableTok}[1]{\textcolor[rgb]{0.07,0.07,0.07}{#1}}
\newcommand{\VerbatimStringTok}[1]{\textcolor[rgb]{0.13,0.47,0.30}{#1}}
\newcommand{\WarningTok}[1]{\textcolor[rgb]{0.37,0.37,0.37}{\textit{#1}}}

\providecommand{\tightlist}{%
  \setlength{\itemsep}{0pt}\setlength{\parskip}{0pt}}
\usepackage{longtable,booktabs,array}
\newcounter{none} % for unnumbered tables
\usepackage{calc} % for calculating minipage widths
% Correct order of tables after \paragraph or \subparagraph
\usepackage{etoolbox}
\makeatletter
\patchcmd\longtable{\par}{\if@noskipsec\mbox{}\fi\par}{}{}
\makeatother
% Allow footnotes in longtable head/foot
\IfFileExists{footnotehyper.sty}{\usepackage{footnotehyper}}{\usepackage{footnote}}
\makesavenoteenv{longtable}

\usepackage{graphicx}
\makeatletter
\newsavebox\pandoc@box
\newcommand*\pandocbounded[1]{% scales image to fit in text height/width
  \sbox\pandoc@box{#1}%
  \Gscale@div\@tempa{\textheight}{\dimexpr\ht\pandoc@box+\dp\pandoc@box\relax}%
  \Gscale@div\@tempb{\linewidth}{\wd\pandoc@box}%
  \ifdim\@tempb\p@<\@tempa\p@\let\@tempa\@tempb\fi% select the smaller of both
  \ifdim\@tempa\p@<\p@\scalebox{\@tempa}{\usebox\pandoc@box}%
  \else\usebox{\pandoc@box}%
  \fi%
}
% Set default figure placement to htbp
\def\fps@figure{htbp}
\makeatother







\usepackage{newtx}

\defaultfontfeatures{Scale=MatchLowercase}
\defaultfontfeatures[\rmfamily]{Ligatures=TeX,Scale=1}





\title{Guía Completa de Vim y Neovim}


\shorttitle{Guía Completa de Vim y Neovim}


\usepackage{etoolbox}



\ccoppy{\textcopyright~2022}



\author{Edison Achalma}



\affiliation{
{Escuela Profesional de Economía, Universidad Nacional de San Cristóbal
de Huamanga}}




\leftheader{Achalma}

\date{2022-09-27}


\abstract{Esta guía exhaustiva presenta una introducción progresiva y
práctica al editor de texto Vim y su evolución moderna, Neovim. Se
abordan desde los fundamentos filosóficos y modales del editor hasta
técnicas avanzadas de edición, navegación eficiente, composición de
comandos, uso de registros, macros, múltiples ventanas y buffers,
búsqueda/reemplazo con expresiones regulares, plegado, marcas y
configuración personalizada tanto en Vimscript como en Lua. Se detallan
las diferencias clave entre ambas versiones, con énfasis en las ventajas
contemporáneas de Neovim; soporte nativo para LSP, Tree-sitter,
arquitectura asíncrona y configuración en Lua. La obra incluye
instrucciones detalladas de instalación en múltiples sistemas
operativos, una selección comentada de plugins esenciales, mapeos
prácticos, trucos avanzados y una hoja de ruta de aprendizaje por
niveles. Dirigida tanto a principiantes que buscan superar la curva
inicial de aprendizaje como a usuarios intermedios-avanzados que desean
optimizar su flujo de trabajo en entornos Linux/Unix, esta guía busca
servir como referencia completa y actualizada para dominar uno de los
editores de texto más potentes y ubicuos del ecosistema del software
libre. }

\keywords{Vim, Neovim, Comandos Vim}

\authornote{\par{\addORCIDlink{Edison Achalma}{0000-0001-6996-3364}} 
\par{ }
\par{   El autor no tiene conflictos de interés que revelar.    Los
roles de autor se clasificaron utilizando la taxonomía de roles de
colaborador (CRediT; https://credit.niso.org/) de la siguiente
manera:  Edison Achalma:   conceptualización, metodología, análisis
formal, investigación, recursos, curación de
datos, redacción, visualización, supervisión, administración del
proyecto}
\par{La correspondencia relativa a este artículo debe dirigirse a Edison
Achalma, Escuela Profesional de Economía, Universidad Nacional de San
Cristóbal de Huamanga, Portal Independencia N°
57, Ayacucho, AYA 5001, Perú, Email: \href{mailto:elmer.achalma.09@unsch.edu.pe}{elmer.achalma.09@unsch.edu.pe}}
}

\makeatletter
\let\endoldlt\endlongtable
\def\endlongtable{
\hline
\endoldlt
}
\makeatother

\urlstyle{same}



\hypersetup{
  pdftitle={Gu\'{\i}a Completa de Vim y Neovim},
  pdfauthor={Autor}
}
\makeatletter
\@ifpackageloaded{caption}{}{\usepackage{caption}}
\AtBeginDocument{%
\ifdefined\contentsname
  \renewcommand*\contentsname{Tabla de contenidos}
\else
  \newcommand\contentsname{Tabla de contenidos}
\fi
\ifdefined\listfigurename
  \renewcommand*\listfigurename{Lista de Figuras}
\else
  \newcommand\listfigurename{Lista de Figuras}
\fi
\ifdefined\listtablename
  \renewcommand*\listtablename{Lista de Tablas}
\else
  \newcommand\listtablename{Lista de Tablas}
\fi
\ifdefined\figurename
  \renewcommand*\figurename{Figura}
\else
  \newcommand\figurename{Figura}
\fi
\ifdefined\tablename
  \renewcommand*\tablename{Tabla}
\else
  \newcommand\tablename{Tabla}
\fi
}
\@ifpackageloaded{float}{}{\usepackage{float}}
\floatstyle{ruled}
\@ifundefined{c@chapter}{\newfloat{codelisting}{h}{lop}}{\newfloat{codelisting}{h}{lop}[chapter]}
\floatname{codelisting}{Listado}
\newcommand*\listoflistings{\listof{codelisting}{Lista de Listados}}
\makeatother
\makeatletter
\makeatother
\makeatletter
\@ifpackageloaded{caption}{}{\usepackage{caption}}
\@ifpackageloaded{subcaption}{}{\usepackage{subcaption}}
\makeatother
\makeatletter
\@ifpackageloaded{fontawesome5}{}{\usepackage{fontawesome5}}
\makeatother

% From https://tex.stackexchange.com/a/645996/211326
%%% apa7 doesn't want to add appendix section titles in the toc
%%% let's make it do it
\makeatletter
\xpatchcmd{\appendix}
  {\par}
  {\addcontentsline{toc}{section}{\@currentlabelname}\par}
  {}{}
\makeatother

%% Disable longtable counter
%% https://tex.stackexchange.com/a/248395/211326

\usepackage{etoolbox}

\makeatletter
\patchcmd{\LT@caption}
  {\bgroup}
  {\bgroup\global\LTpatch@captiontrue}
  {}{}
\patchcmd{\longtable}
  {\par}
  {\par\global\LTpatch@captionfalse}
  {}{}
\apptocmd{\endlongtable}
  {\ifLTpatch@caption\else\addtocounter{table}{-1}\fi}
  {}{}
\newif\ifLTpatch@caption
\makeatother

\begin{document}

\maketitle


\hypertarget{toc}{}
\tableofcontents
\newpage
\section[Introduction]{Guía Completa de Vim y Neovim}

\setcounter{secnumdepth}{3}

\setlength\LTleft{0pt}




\textbf{¿Qué es Vim?}

Vim (Vi IMproved) es un editor de texto modal altamente configurable,
creado por Bram Moolenaar como una versión mejorada del editor Unix
\texttt{vi}. Es conocido por:

\begin{itemize}
\tightlist
\item
  \textbf{Eficiencia}: Una vez dominado, permite editar texto
  extremadamente rápido
\item
  \textbf{Ubicuidad}: Disponible en prácticamente todos los sistemas
  Unix/Linux
\item
  \textbf{Potencia}: Capacidades avanzadas de edición y automatización
\item
  \textbf{Extensibilidad}: Sistema robusto de plugins y configuración
\end{itemize}

\textbf{¿Qué es Neovim?}

Neovim es una refactorización moderna de Vim que mantiene compatibilidad
mientras añade:

\begin{itemize}
\tightlist
\item
  \textbf{Arquitectura moderna}: Cliente-servidor, API asíncrona
\item
  \textbf{LSP integrado}: Language Server Protocol nativo
\item
  \textbf{Lua}: Lenguaje de configuración más rápido que Vimscript
\item
  \textbf{UI remota}: Soporte para interfaces gráficas externas
\item
  \textbf{Rendimiento}: Más rápido y eficiente
\item
  \textbf{Desarrollo activo}: Comunidad muy activa
\end{itemize}

\textbf{Diferencias Principales}

{\def\LTcaptype{none} % do not increment counter
\begin{longtable}[]{@{}lll@{}}
\toprule\noalign{}
Característica & Vim & Neovim \\
\midrule\noalign{}
\endhead
\bottomrule\noalign{}
\endlastfoot
Lenguaje config & Vimscript & Vimscript + Lua \\
LSP & Plugin & Integrado \\
Async & Plugin & Nativo \\
Tree-sitter & No & Sí \\
UI externa & Limitada & Completa \\
Desarrollo & Conservador & Agresivo \\
\end{longtable}
}

\section{Instalación}\label{instalaciuxf3n}

\subsection{Instalar Vim}\label{instalar-vim}

\textbf{Ubuntu/Debian:}

\begin{Shaded}
\begin{Highlighting}[]
\FunctionTok{sudo}\NormalTok{ apt update}
\FunctionTok{sudo}\NormalTok{ apt install vim}
\end{Highlighting}
\end{Shaded}

\textbf{Fedora:}

\begin{Shaded}
\begin{Highlighting}[]
\FunctionTok{sudo}\NormalTok{ dnf install vim{-}enhanced}
\end{Highlighting}
\end{Shaded}

\textbf{macOS:}

\begin{Shaded}
\begin{Highlighting}[]
\ExtensionTok{brew}\NormalTok{ install vim}
\end{Highlighting}
\end{Shaded}

\textbf{Windows:} Descargar desde: https://www.vim.org/download.php

\subsection{Instalar Neovim}\label{instalar-neovim}

\textbf{Ubuntu/Debian (PPA):}

\begin{Shaded}
\begin{Highlighting}[]
\FunctionTok{sudo}\NormalTok{ add{-}apt{-}repository ppa:neovim{-}ppa/stable}
\FunctionTok{sudo}\NormalTok{ apt update}
\FunctionTok{sudo}\NormalTok{ apt install neovim}
\end{Highlighting}
\end{Shaded}

\textbf{Fedora:}

\begin{Shaded}
\begin{Highlighting}[]
\FunctionTok{sudo}\NormalTok{ dnf install neovim}
\end{Highlighting}
\end{Shaded}

\textbf{macOS:}

\begin{Shaded}
\begin{Highlighting}[]
\ExtensionTok{brew}\NormalTok{ install neovim}
\end{Highlighting}
\end{Shaded}

\textbf{Arch Linux:}

\begin{Shaded}
\begin{Highlighting}[]
\FunctionTok{sudo}\NormalTok{ pacman }\AttributeTok{{-}S}\NormalTok{ neovim}
\end{Highlighting}
\end{Shaded}

\textbf{Desde AppImage (cualquier Linux):}

\begin{Shaded}
\begin{Highlighting}[]
\ExtensionTok{curl} \AttributeTok{{-}LO}\NormalTok{ https://github.com/neovim/neovim/releases/latest/download/nvim.appimage}
\FunctionTok{chmod}\NormalTok{ u+x nvim.appimage}
\ExtensionTok{./nvim.appimage}
\end{Highlighting}
\end{Shaded}

\textbf{Windows:}

\begin{Shaded}
\begin{Highlighting}[]
\ExtensionTok{choco}\NormalTok{ install neovim}
\CommentTok{\# O}
\ExtensionTok{scoop}\NormalTok{ install neovim}
\end{Highlighting}
\end{Shaded}

\subsection{Verificar Instalación}\label{verificar-instalaciuxf3n}

\begin{Shaded}
\begin{Highlighting}[]
\CommentTok{\# Vim}
\ExtensionTok{vim} \AttributeTok{{-}{-}version}

\CommentTok{\# Neovim}
\ExtensionTok{nvim} \AttributeTok{{-}{-}version}
\end{Highlighting}
\end{Shaded}

\subsection{Configurar Alias}\label{configurar-alias}

Para usar \texttt{vim} para lanzar Neovim:

\begin{Shaded}
\begin{Highlighting}[]
\CommentTok{\# En \textasciitilde{}/.bashrc o \textasciitilde{}/.zshrc}
\BuiltInTok{alias}\NormalTok{ vim=}\StringTok{\textquotesingle{}nvim\textquotesingle{}}
\BuiltInTok{alias}\NormalTok{ vi=}\StringTok{\textquotesingle{}nvim\textquotesingle{}}
\end{Highlighting}
\end{Shaded}

\section{Filosofía y Conceptos
Básicos}\label{filosofuxeda-y-conceptos-buxe1sicos}

\subsection{Filosofía Modal}\label{filosofuxeda-modal}

Vim usa \textbf{modos} diferentes para distintas tareas:

\begin{itemize}
\tightlist
\item
  \textbf{Normal}: Navegar y ejecutar comandos (modo por defecto)
\item
  \textbf{Insert}: Insertar texto
\item
  \textbf{Visual}: Seleccionar texto
\item
  \textbf{Command}: Ejecutar comandos Ex
\end{itemize}

Esta separación permite:

\begin{itemize}
\tightlist
\item
  \textbf{Eficiencia}: Cada tecla hace algo útil en modo Normal
\item
  \textbf{Composición}: Combinar operadores con movimientos
\item
  \textbf{Repetición}: El comando \texttt{.} repite la última acción
\end{itemize}

\subsection{Composición de Comandos}\label{composiciuxf3n-de-comandos}

Vim usa una gramática:

\begin{verbatim}
[operador][movimiento]
[operador][número][movimiento]
\end{verbatim}

\textbf{Ejemplos:}

\begin{itemize}
\tightlist
\item
  \texttt{d2w} = Eliminar (delete) 2 palabras (words)
\item
  \texttt{c\$} = Cambiar (change) hasta fin de línea
\item
  \texttt{y3j} = Copiar (yank) 3 líneas hacia abajo
\end{itemize}

\subsection{Filosofía ``Modal''}\label{filosofuxeda-modal-1}

\begin{itemize}
\tightlist
\item
  \textbf{Separación de tareas}: No mezclar edición y navegación
\item
  \textbf{Comando vs texto}: Modo Normal para comandos, Insert para
  texto
\item
  \textbf{Eficiencia sobre simplicidad}: Curva de aprendizaje pero
  altamente eficiente
\end{itemize}

\subsection{El Camino del Vim}\label{el-camino-del-vim}

\begin{enumerate}
\def\labelenumi{\arabic{enumi}.}
\tightlist
\item
  \textbf{Semana 1-2}: Básicos (hjkl, i, a, o, Esc, :wq)
\item
  \textbf{Mes 1}: Movimientos (w, b, f, t, /, n)
\item
  \textbf{Meses 2-3}: Operadores (d, c, y con movimientos)
\item
  \textbf{Meses 4-6}: Visual, macros, registros
\item
  \textbf{Después}: Plugins, configuración avanzada, Vimscript/Lua
\end{enumerate}

\section{Modos de Vim}\label{modos-de-vim}

\subsection{Modo Normal (Default)}\label{modo-normal-default}

El modo por defecto. Todas las teclas son comandos. Para activar el modo
comando en Vim, debes presionar la tecla ``Esc''. Esto te llevará al
modo comando desde cualquier otro modo en el que te encuentres, como el
modo insertar o el modo de reemplazo. Una vez que estés en el modo
comando, puedes utilizar una variedad de comandos y combinaciones de
teclas para navegar, editar y guardar tus archivos. Para salir de Vim,
puedes ingresar el comando ``:q'' seguido de Enter. Si has realizado
cambios y deseas guardarlos antes de salir, utiliza el comando ``:wq''
para escribir y guardar los cambios y salir de Vim.

\textbf{Entrar:}

\begin{itemize}
\tightlist
\item
  \texttt{Esc} desde cualquier modo
\item
  \texttt{Ctrl+{[}} (alternativa a Esc)
\item
  \texttt{Ctrl+c} (similar a Esc)
\end{itemize}

\subsection{Modo Insert}\label{modo-insert}

Para insertar texto.

\textbf{Entrar desde Normal:}

\begin{itemize}
\tightlist
\item
  \texttt{i} - Insertar antes del cursor
\item
  \texttt{I} - Insertar al inicio de la línea
\item
  \texttt{a} - Insertar después del cursor (append)
\item
  \texttt{A} - Insertar al final de la línea
\item
  \texttt{o} - Abrir línea nueva abajo
\item
  \texttt{O} - Abrir línea nueva arriba
\item
  \texttt{s} - Sustituir carácter (eliminar + insert)
\item
  \texttt{S} - Sustituir línea completa
\item
  \texttt{C} - Cambiar hasta fin de línea (c\$)
\end{itemize}

\textbf{Salir:}

\begin{itemize}
\tightlist
\item
  \texttt{Esc} o \texttt{Ctrl+{[}} - Volver a Normal
\end{itemize}

\subsection{Modo Visual}\label{modo-visual}

Para seleccionar texto.

\textbf{Entrar desde Normal:}

\begin{itemize}
\tightlist
\item
  \texttt{v} - Visual por carácter
\item
  \texttt{V} - Visual por línea
\item
  \texttt{Ctrl+v} - Visual en bloque
\end{itemize}

\textbf{Operaciones:}

Una vez seleccionado:

\begin{itemize}
\tightlist
\item
  \texttt{d} - Eliminar
\item
  \texttt{c} - Cambiar
\item
  \texttt{y} - Copiar
\item
  \texttt{\textgreater{}} - Indentar
\item
  \texttt{\textless{}} - Desindentar
\item
  \texttt{\textasciitilde{}} - Cambiar mayúsculas/minúsculas
\end{itemize}

\subsection{Modo Command
(Command-line)}\label{modo-command-command-line}

Para ejecutar comandos Ex.

\textbf{Entrar desde Normal:}

\begin{itemize}
\tightlist
\item
  \texttt{:} - Comando Ex
\item
  \texttt{/} - Búsqueda hacia adelante
\item
  \texttt{?} - Búsqueda hacia atrás
\end{itemize}

\textbf{Comandos comunes:}

\begin{itemize}
\tightlist
\item
  \texttt{:w} - Guardar (write)
\item
  \texttt{:q} - Salir (quit)
\item
  \texttt{:wq} o \texttt{:x} - Guardar y salir
\item
  \texttt{:q!} - Salir sin guardar
\item
  \texttt{:e\ archivo} - Abrir archivo (edit)
\item
  \texttt{:help\ tema} - Ayuda
\end{itemize}

\subsection{Modo Replace}\label{modo-replace}

Sobrescribir texto existente.

\textbf{Entrar desde Normal:}

\begin{itemize}
\tightlist
\item
  \texttt{R} - Modo replace (sobrescribir)
\item
  \texttt{r} - Reemplazar un solo carácter
\end{itemize}

\subsection{Modo Select}\label{modo-select}

Similar a selección en editores normales.

\textbf{Entrar:}

\begin{itemize}
\tightlist
\item
  \texttt{gh} - Select mode por carácter
\item
  \texttt{gH} - Select mode por línea
\item
  \texttt{g\ Ctrl+h} - Select mode en bloque
\end{itemize}

\section{Navegación Básica}\label{navegaciuxf3n-buxe1sica}

\subsection{Movimiento de Cursor (Modo
Normal)}\label{movimiento-de-cursor-modo-normal}

\subsection{Básico (hjkl)}\label{buxe1sico-hjkl}

{\def\LTcaptype{none} % do not increment counter
\begin{longtable}[]{@{}ll@{}}
\toprule\noalign{}
Tecla & Movimiento \\
\midrule\noalign{}
\endhead
\bottomrule\noalign{}
\endlastfoot
\texttt{h} & ← Izquierda \\
\texttt{j} & ↓ Abajo \\
\texttt{k} & ↑ Arriba \\
\texttt{l} & → Derecha \\
\end{longtable}
}

\subsection{Por Palabras}\label{por-palabras}

{\def\LTcaptype{none} % do not increment counter
\begin{longtable}[]{@{}ll@{}}
\toprule\noalign{}
Tecla & Movimiento \\
\midrule\noalign{}
\endhead
\bottomrule\noalign{}
\endlastfoot
\texttt{w} & Inicio de siguiente palabra \\
\texttt{W} & Inicio de siguiente PALABRA (ignora puntuación) \\
\texttt{e} & Final de palabra actual/siguiente \\
\texttt{E} & Final de PALABRA \\
\texttt{b} & Inicio de palabra anterior \\
\texttt{B} & Inicio de PALABRA anterior \\
\texttt{ge} & Final de palabra anterior \\
\end{longtable}
}

\textbf{Nota:} ``palabra'' = letras/números, ``PALABRA'' = cualquier
carácter no-espacio

\subsection{Por Línea}\label{por-luxednea}

{\def\LTcaptype{none} % do not increment counter
\begin{longtable}[]{@{}ll@{}}
\toprule\noalign{}
Tecla & Movimiento \\
\midrule\noalign{}
\endhead
\bottomrule\noalign{}
\endlastfoot
\texttt{0} & Inicio de línea (columna 0) \\
\texttt{\^{}} & Primer carácter no-blanco \\
\texttt{\$} & Final de línea \\
\texttt{g\_} & Último carácter no-blanco \\
\texttt{g0} & Inicio de línea visual (con wrap) \\
\texttt{g\$} & Final de línea visual \\
\end{longtable}
}

\subsection{Por Pantalla}\label{por-pantalla}

{\def\LTcaptype{none} % do not increment counter
\begin{longtable}[]{@{}ll@{}}
\toprule\noalign{}
Tecla & Movimiento \\
\midrule\noalign{}
\endhead
\bottomrule\noalign{}
\endlastfoot
\texttt{H} & Top de pantalla (High) \\
\texttt{M} & Medio de pantalla (Middle) \\
\texttt{L} & Bottom de pantalla (Low) \\
\texttt{zt} & Scroll: línea actual arriba \\
\texttt{zz} & Scroll: línea actual al centro \\
\texttt{zb} & Scroll: línea actual abajo \\
\texttt{Ctrl+e} & Scroll abajo una línea \\
\texttt{Ctrl+y} & Scroll arriba una línea \\
\texttt{Ctrl+d} & Media página abajo \\
\texttt{Ctrl+u} & Media página arriba \\
\texttt{Ctrl+f} & Página completa abajo (forward) \\
\texttt{Ctrl+b} & Página completa arriba (backward) \\
\end{longtable}
}

\subsection{Por Archivo}\label{por-archivo}

{\def\LTcaptype{none} % do not increment counter
\begin{longtable}[]{@{}ll@{}}
\toprule\noalign{}
Tecla & Movimiento \\
\midrule\noalign{}
\endhead
\bottomrule\noalign{}
\endlastfoot
\texttt{gg} & Inicio del archivo \\
\texttt{G} & Final del archivo \\
\texttt{:\textless{}n\textgreater{}} o
\texttt{\textless{}n\textgreater{}G} & Ir a línea n \\
\texttt{\%} & Ir al paréntesis/llave/corchete coincidente \\
\end{longtable}
}

\subsection{Búsqueda de Caracteres en
Línea}\label{buxfasqueda-de-caracteres-en-luxednea}

{\def\LTcaptype{none} % do not increment counter
\begin{longtable}[]{@{}ll@{}}
\toprule\noalign{}
Tecla & Movimiento \\
\midrule\noalign{}
\endhead
\bottomrule\noalign{}
\endlastfoot
\texttt{f\{char\}} & Find: ir a siguiente \{char\} \\
\texttt{F\{char\}} & Find: ir a anterior \{char\} \\
\texttt{t\{char\}} & Till: ir antes de siguiente \{char\} \\
\texttt{T\{char\}} & Till: ir después de anterior \{char\} \\
\texttt{;} & Repetir último f/F/t/T \\
\texttt{,} & Repetir último f/F/t/T en reversa \\
\end{longtable}
}

\textbf{Ejemplos:}

\begin{verbatim}
Esta es una línea de ejemplo
\end{verbatim}

Con cursor en `E':

\begin{itemize}
\tightlist
\item
  \texttt{fw} - Ir a `u' (primera `u')
\item
  \texttt{3fl} - Ir a tercera `l'
\item
  \texttt{t.} - Ir antes del punto
\end{itemize}

\section{Edición de Texto}\label{ediciuxf3n-de-texto}

\subsection{Insertar}\label{insertar}

{\def\LTcaptype{none} % do not increment counter
\begin{longtable}[]{@{}ll@{}}
\toprule\noalign{}
Comando & Acción \\
\midrule\noalign{}
\endhead
\bottomrule\noalign{}
\endlastfoot
\texttt{i} & Insertar antes del cursor \\
\texttt{I} & Insertar al inicio de línea \\
\texttt{a} & Insertar después del cursor (append) \\
\texttt{A} & Insertar al final de línea \\
\texttt{o} & Abrir línea nueva abajo \\
\texttt{O} & Abrir línea nueva arriba \\
\texttt{gi} & Insertar en última posición de inserción \\
\end{longtable}
}

\subsection{Eliminar}\label{eliminar}

{\def\LTcaptype{none} % do not increment counter
\begin{longtable}[]{@{}ll@{}}
\toprule\noalign{}
Comando & Acción \\
\midrule\noalign{}
\endhead
\bottomrule\noalign{}
\endlastfoot
\texttt{x} & Eliminar carácter bajo cursor \\
\texttt{X} & Eliminar carácter antes del cursor \\
\texttt{dd} & Eliminar línea \\
\texttt{D} & Eliminar hasta fin de línea \texttt{d\$} \\
\texttt{d\{movimiento\}} & Eliminar hasta movimiento \\
\texttt{dw} & Eliminar palabra \\
\texttt{db} & Eliminar palabra hacia atrás \\
\texttt{d\$} o \texttt{D} & Eliminar hasta fin de línea \\
\texttt{d0} & Eliminar hasta inicio de línea \\
\texttt{dG} & Eliminar hasta fin de archivo \\
\texttt{dgg} & Eliminar hasta inicio de archivo \\
\end{longtable}
}

\subsection{Cambiar (Change)}\label{cambiar-change}

Igual que eliminar pero entra en modo Insert.

{\def\LTcaptype{none} % do not increment counter
\begin{longtable}[]{@{}ll@{}}
\toprule\noalign{}
Comando & Acción \\
\midrule\noalign{}
\endhead
\bottomrule\noalign{}
\endlastfoot
\texttt{c\{movimiento\}} & Cambiar texto \\
\texttt{cc} o \texttt{S} & Cambiar línea completa \\
\texttt{C} & Cambiar hasta fin de línea \\
\texttt{cw} & Cambiar palabra \\
\texttt{ciw} & Cambiar palabra interna (inner word) \\
\texttt{ci"} & Cambiar dentro de comillas \\
\texttt{ci(} & Cambiar dentro de paréntesis \\
\texttt{cit} & Cambiar dentro de tag HTML \\
\end{longtable}
}

\subsection{Copiar y Pegar (Yank/Put)}\label{copiar-y-pegar-yankput}

{\def\LTcaptype{none} % do not increment counter
\begin{longtable}[]{@{}ll@{}}
\toprule\noalign{}
Comando & Acción \\
\midrule\noalign{}
\endhead
\bottomrule\noalign{}
\endlastfoot
\texttt{y\{movimiento\}} & Copiar (yank) \\
\texttt{yy} o \texttt{Y} & Copiar línea \\
\texttt{yw} & Copiar palabra \\
\texttt{y\$} & Copiar hasta fin de línea \\
\texttt{p} & Pegar después del cursor \\
\texttt{P} & Pegar antes del cursor \\
\texttt{gp} & Pegar y mover cursor después \\
\texttt{gP} & Pegar antes y mover cursor \\
\end{longtable}
}

\subsection{Reemplazar}\label{reemplazar}

{\def\LTcaptype{none} % do not increment counter
\begin{longtable}[]{@{}ll@{}}
\toprule\noalign{}
Comando & Acción \\
\midrule\noalign{}
\endhead
\bottomrule\noalign{}
\endlastfoot
\texttt{r\{char\}} & Reemplazar carácter bajo cursor \\
\texttt{R} & Entrar en modo Replace \\
\texttt{\textasciitilde{}} & Cambiar mayúscula/minúscula \\
\texttt{g\textasciitilde{}\{movimiento\}} & Cambiar caso \\
\texttt{gu\{movimiento\}} & A minúsculas \\
\texttt{gU\{movimiento\}} & A MAYÚSCULAS \\
\end{longtable}
}

\subsection{Deshacer y Rehacer}\label{deshacer-y-rehacer}

{\def\LTcaptype{none} % do not increment counter
\begin{longtable}[]{@{}ll@{}}
\toprule\noalign{}
Comando & Acción \\
\midrule\noalign{}
\endhead
\bottomrule\noalign{}
\endlastfoot
\texttt{u} & Deshacer (undo) \\
\texttt{Ctrl+r} & Rehacer (redo) \\
\texttt{U} & Deshacer todos los cambios en línea \\
\end{longtable}
}

\subsection{Repetir}\label{repetir}

{\def\LTcaptype{none} % do not increment counter
\begin{longtable}[]{@{}ll@{}}
\toprule\noalign{}
Comando & Acción \\
\midrule\noalign{}
\endhead
\bottomrule\noalign{}
\endlastfoot
\texttt{.} & Repetir último cambio \\
\texttt{@:} & Repetir último comando Ex \\
\texttt{@@} & Repetir último macro \\
\end{longtable}
}

\subsection{Unir Líneas}\label{unir-luxedneas}

{\def\LTcaptype{none} % do not increment counter
\begin{longtable}[]{@{}ll@{}}
\toprule\noalign{}
Comando & Acción \\
\midrule\noalign{}
\endhead
\bottomrule\noalign{}
\endlastfoot
\texttt{J} & Unir línea siguiente con actual \\
\texttt{gJ} & Unir sin añadir espacio \\
\end{longtable}
}

\subsection{Indentación}\label{indentaciuxf3n}

{\def\LTcaptype{none} % do not increment counter
\begin{longtable}[]{@{}ll@{}}
\toprule\noalign{}
Comando & Acción \\
\midrule\noalign{}
\endhead
\bottomrule\noalign{}
\endlastfoot
\texttt{\textgreater{}\textgreater{}} & Indentar línea \\
\texttt{\textless{}\textless{}} & Desindentar línea \\
\texttt{==} & Auto-indentar línea \\
\texttt{\textgreater{}\%} & Indentar bloque (con cursor en
paréntesis) \\
\texttt{gg=G} & Auto-indentar archivo completo \\
\end{longtable}
}

\section{Comandos de Búsqueda}\label{comandos-de-buxfasqueda}

\subsection{Búsqueda Básica}\label{buxfasqueda-buxe1sica}

{\def\LTcaptype{none} % do not increment counter
\begin{longtable}[]{@{}ll@{}}
\toprule\noalign{}
Comando & Acción \\
\midrule\noalign{}
\endhead
\bottomrule\noalign{}
\endlastfoot
\texttt{/patrón} & Buscar hacia adelante \\
\texttt{?patrón} & Buscar hacia atrás \\
\texttt{n} & Siguiente ocurrencia \\
\texttt{N} & Ocurrencia anterior \\
\texttt{*} & Buscar palabra bajo cursor (adelante) \\
\texttt{\#} & Buscar palabra bajo cursor (atrás) \\
\texttt{g*} & Buscar parcial hacia adelante \\
\texttt{g\#} & Buscar parcial hacia atrás \\
\end{longtable}
}

\subsection{Opciones de Búsqueda}\label{opciones-de-buxfasqueda}

\begin{Shaded}
\begin{Highlighting}[]
\NormalTok{:set ignorecase   " Ignorar mayúsculas/minúsculas}
\NormalTok{:set smartcase    " Case{-}sensitive si hay mayúsculas}
\NormalTok{:set incsearch    " Búsqueda incremental}
\NormalTok{:set hlsearch     " Resaltar búsquedas}
\NormalTok{:nohlsearch       " Quitar resaltado (o :noh)}
\end{Highlighting}
\end{Shaded}

\subsection{Expresiones Regulares}\label{expresiones-regulares}

\textbf{Metacaracteres básicos:}

\begin{itemize}
\tightlist
\item
  \texttt{.} - Cualquier carácter
\item
  \texttt{*} - 0 o más del anterior
\item
  \texttt{\^{}} - Inicio de línea
\item
  \texttt{\$} - Fin de línea
\item
  \texttt{{[}{]}} - Conjunto de caracteres
\item
  \texttt{\textbackslash{}\textless{}} - Inicio de palabra
\item
  \texttt{\textbackslash{}\textgreater{}} - Fin de palabra
\end{itemize}

\textbf{Ejemplos:}

\begin{Shaded}
\begin{Highlighting}[]
\NormalTok{/\textbackslash{}\textless{}vim\textbackslash{}\textgreater{}        " Palabra exacta "vim"}
\NormalTok{/\^{}\#             " Líneas que empiezan con \#}
\NormalTok{/error.*$       " "error" hasta fin de línea}
\NormalTok{/[0{-}9]\textbackslash{}+        " Uno o más dígitos}
\NormalTok{/\textbackslash{}d\textbackslash{}\{3\}{-}\textbackslash{}d\textbackslash{}\{4\}  " XXX{-}XXXX (teléfono)}
\end{Highlighting}
\end{Shaded}

\section{Modo Visual}\label{modo-visual-1}

\subsection{Tipos de Selección
Visual}\label{tipos-de-selecciuxf3n-visual}

{\def\LTcaptype{none} % do not increment counter
\begin{longtable}[]{@{}ll@{}}
\toprule\noalign{}
Comando & Tipo \\
\midrule\noalign{}
\endhead
\bottomrule\noalign{}
\endlastfoot
\texttt{v} & Visual por carácter \\
\texttt{V} & Visual por línea \\
\texttt{Ctrl+v} & Visual en bloque (columnar) \\
\texttt{gv} & Reseleccionar última selección \\
\end{longtable}
}

\subsection{Operaciones en Visual}\label{operaciones-en-visual}

Una vez en modo visual:

{\def\LTcaptype{none} % do not increment counter
\begin{longtable}[]{@{}ll@{}}
\toprule\noalign{}
Comando & Acción \\
\midrule\noalign{}
\endhead
\bottomrule\noalign{}
\endlastfoot
\texttt{d} & Eliminar selección \\
\texttt{c} & Cambiar selección \\
\texttt{y} & Copiar selección \\
\texttt{\textgreater{}} & Indentar \\
\texttt{\textless{}} & Desindentar \\
\texttt{=} & Auto-indentar \\
\texttt{\textasciitilde{}} & Cambiar mayúsculas/minúsculas \\
\texttt{u} & A minúsculas \\
\texttt{U} & A MAYÚSCULAS \\
\texttt{o} & Ir al otro extremo de selección \\
\texttt{O} & Ir a otra esquina (bloque) \\
\end{longtable}
}

\subsection{Visual Block (Columnar)}\label{visual-block-columnar}

\textbf{Casos de uso:}

\begin{enumerate}
\def\labelenumi{\arabic{enumi}.}
\tightlist
\item
  \textbf{Insertar en múltiples líneas}
\item
  \textbf{Eliminar columnas}
\item
  \textbf{Cambiar bloques}
\end{enumerate}

\textbf{Ejemplo - Insertar en múltiples líneas:}

\begin{verbatim}
1. Ctrl+v (visual block)
2. Seleccionar líneas (j/k)
3. I (Insert al inicio)
4. Escribir texto
5. Esc (se aplica a todas las líneas)
\end{verbatim}

\subsection{Objetos de Texto}\label{objetos-de-texto}

Operan sobre ``objetos'' como palabras, párrafos, etc.

\textbf{Sintaxis:} \texttt{\{operador\}\{a/i\}\{objeto\}}

{\def\LTcaptype{none} % do not increment counter
\begin{longtable}[]{@{}ll@{}}
\toprule\noalign{}
Objeto & Descripción \\
\midrule\noalign{}
\endhead
\bottomrule\noalign{}
\endlastfoot
\texttt{w} & Palabra (word) \\
\texttt{W} & PALABRA \\
\texttt{s} & Sentencia \\
\texttt{p} & Párrafo \\
\texttt{{[}} o \texttt{{]}} & Bloque {[}{]} \\
\texttt{(} o \texttt{)} & Bloque () \\
\texttt{\{} o \texttt{\}} & Bloque \{\} \\
\texttt{\textless{}} o \texttt{\textgreater{}} & Bloque
\textless\textgreater{} \\
\texttt{t} & Tag HTML \\
\texttt{"} & Comillas dobles \\
\texttt{\textquotesingle{}} & Comillas simples \\
\texttt{\textasciigrave{}} & Backticks \\
\end{longtable}
}

\textbf{Modificadores:}

\begin{itemize}
\tightlist
\item
  \texttt{a} - ``a'' (around) - incluye delimitadores
\item
  \texttt{i} - ``inner'' - solo contenido
\end{itemize}

\textbf{Ejemplos:}

\begin{Shaded}
\begin{Highlighting}[]
\NormalTok{ciw    " Change inner word}
\NormalTok{da"    " Delete a quote (incluye comillas)}
\NormalTok{yi(    " Yank inner parentheses}
\NormalTok{vit    " Visual select inner tag}
\NormalTok{ci\{    " Change inner braces}
\end{Highlighting}
\end{Shaded}

\section{Registros y Macros}\label{registros-y-macros}

\subsection{Registros}\label{registros}

Vim tiene múltiples ``clipboards'' llamados registros.

\subsection{Tipos de Registros}\label{tipos-de-registros}

{\def\LTcaptype{none} % do not increment counter
\begin{longtable}[]{@{}ll@{}}
\toprule\noalign{}
Registro & Descripción \\
\midrule\noalign{}
\endhead
\bottomrule\noalign{}
\endlastfoot
\texttt{""} & Sin nombre (defecto) \\
\texttt{"0} & Último yank \\
\texttt{"1}-\texttt{"9} & Historial de deletes \\
\texttt{"a}-\texttt{"z} & Nombrados (minúsculas) \\
\texttt{"A}-\texttt{"Z} & Nombrados (append) \\
\texttt{"\_} & Agujero negro (no guarda) \\
\texttt{"+} & Clipboard del sistema \\
\texttt{"*} & Clipboard de selección (X11) \\
\texttt{"\%} & Nombre del archivo actual \\
\texttt{"\#} & Nombre del archivo alterno \\
\texttt{"/} & Último patrón de búsqueda \\
\texttt{":} & Último comando Ex \\
\end{longtable}
}

\subsection{Usar Registros}\label{usar-registros}

\textbf{Sintaxis:} \texttt{"\{registro\}\{comando\}}

\begin{Shaded}
\begin{Highlighting}[]
\NormalTok{"ayy       " Copiar línea al registro a}
\NormalTok{"ap        " Pegar desde registro a}
\NormalTok{"Ayy       " Añadir línea al registro a}
\NormalTok{"\_dd       " Eliminar sin guardar}
\NormalTok{"+y        " Copiar al clipboard del sistema}
\NormalTok{"+p        " Pegar desde clipboard}
\end{Highlighting}
\end{Shaded}

\subsection{Ver Registros}\label{ver-registros}

\begin{Shaded}
\begin{Highlighting}[]
\NormalTok{:registers      " Ver todos los registros}
\NormalTok{:reg a b c      " Ver registros específicos}
\end{Highlighting}
\end{Shaded}

\subsection{Macros}\label{macros}

Grabar y repetir secuencias de comandos.

\subsection{Grabar Macro}\label{grabar-macro}

\begin{Shaded}
\begin{Highlighting}[]
\NormalTok{q\{registro\}     " Empezar a grabar en registro}
\NormalTok{... comandos ...}
\NormalTok{q               " Terminar grabación}
\end{Highlighting}
\end{Shaded}

\subsection{Ejecutar Macro}\label{ejecutar-macro}

\begin{Shaded}
\begin{Highlighting}[]
\NormalTok{@\{registro\}     " Ejecutar macro}
\NormalTok{@@              " Repetir último macro}
\NormalTok{\{n\}@\{registro\}  " Ejecutar n veces}
\end{Highlighting}
\end{Shaded}

\textbf{Ejemplo - Formatear lista:}

\begin{Shaded}
\begin{Highlighting}[]
\NormalTok{\# Lista:}
\NormalTok{item 1}
\NormalTok{item 2}
\NormalTok{item 3}

\NormalTok{\# Grabar macro en registro a:}
\NormalTok{qa              " Empezar grabación}
\NormalTok{I{-} \textless{}Esc\textgreater{}        " Insertar "{-} " al inicio}
\NormalTok{j               " Bajar una línea}
\NormalTok{q               " Terminar grabación}

\NormalTok{\# Ejecutar:}
\NormalTok{2@a             " Aplicar a siguientes 2 líneas}

\NormalTok{\# Resultado:}
\NormalTok{{-} item 1}
\NormalTok{{-} item 2}
\NormalTok{{-} item 3}
\end{Highlighting}
\end{Shaded}

\subsection{Editar Macros}\label{editar-macros}

\begin{Shaded}
\begin{Highlighting}[]
\NormalTok{\# Ver contenido del registro}
\NormalTok{:echo @a}

\NormalTok{\# Editar}
\NormalTok{:let @a = \textquotesingle{}nuevo contenido\textquotesingle{}}

\NormalTok{\# O pegar en buffer, editar, y copiar de vuelta}
\NormalTok{"ap             " Pegar macro}
\NormalTok{... editar ...}
\NormalTok{"ayy            " Copiar de vuelta}
\end{Highlighting}
\end{Shaded}

\section{Ventanas y Pestañas}\label{ventanas-y-pestauxf1as}

\subsection{Ventanas (Splits)}\label{ventanas-splits}

\subsection{Crear Ventanas}\label{crear-ventanas}

{\def\LTcaptype{none} % do not increment counter
\begin{longtable}[]{@{}ll@{}}
\toprule\noalign{}
Comando & Acción \\
\midrule\noalign{}
\endhead
\bottomrule\noalign{}
\endlastfoot
\texttt{:split} o \texttt{:sp} & Dividir horizontalmente \\
\texttt{:vsplit} o \texttt{:vsp} & Dividir verticalmente \\
\texttt{:new} & Nueva ventana horizontal \\
\texttt{:vnew} & Nueva ventana vertical \\
\texttt{Ctrl+w\ s} & Split horizontal \\
\texttt{Ctrl+w\ v} & Split vertical \\
\texttt{:sp\ archivo} & Abrir archivo en split \\
\end{longtable}
}

\subsection{Navegar Entre Ventanas}\label{navegar-entre-ventanas}

{\def\LTcaptype{none} % do not increment counter
\begin{longtable}[]{@{}ll@{}}
\toprule\noalign{}
Comando & Acción \\
\midrule\noalign{}
\endhead
\bottomrule\noalign{}
\endlastfoot
\texttt{Ctrl+w\ h} & Ir a ventana izquierda \\
\texttt{Ctrl+w\ j} & Ir a ventana abajo \\
\texttt{Ctrl+w\ k} & Ir a ventana arriba \\
\texttt{Ctrl+w\ l} & Ir a ventana derecha \\
\texttt{Ctrl+w\ w} & Ciclar entre ventanas \\
\texttt{Ctrl+w\ p} & Ir a ventana anterior \\
\end{longtable}
}

\subsection{Mover Ventanas}\label{mover-ventanas}

{\def\LTcaptype{none} % do not increment counter
\begin{longtable}[]{@{}ll@{}}
\toprule\noalign{}
Comando & Acción \\
\midrule\noalign{}
\endhead
\bottomrule\noalign{}
\endlastfoot
\texttt{Ctrl+w\ H} & Mover ventana a la izquierda \\
\texttt{Ctrl+w\ J} & Mover ventana abajo \\
\texttt{Ctrl+w\ K} & Mover ventana arriba \\
\texttt{Ctrl+w\ L} & Mover ventana a la derecha \\
\texttt{Ctrl+w\ r} & Rotar ventanas \\
\texttt{Ctrl+w\ x} & Intercambiar con siguiente \\
\end{longtable}
}

\subsection{Redimensionar Ventanas}\label{redimensionar-ventanas}

{\def\LTcaptype{none} % do not increment counter
\begin{longtable}[]{@{}ll@{}}
\toprule\noalign{}
Comando & Acción \\
\midrule\noalign{}
\endhead
\bottomrule\noalign{}
\endlastfoot
\texttt{Ctrl+w\ =} & Igualar tamaños \\
\texttt{Ctrl+w\ +} & Aumentar altura \\
\texttt{Ctrl+w\ -} & Disminuir altura \\
\texttt{Ctrl+w\ \textgreater{}} & Aumentar ancho \\
\texttt{Ctrl+w\ \textless{}} & Disminuir ancho \\
\texttt{Ctrl+w\ \_} & Maximizar altura \\
\texttt{Ctrl+w\ \textbar{}} & Maximizar ancho \\
\texttt{:resize\ \{n\}} & Establecer altura \\
\texttt{:vertical\ resize\ \{n\}} & Establecer ancho \\
\end{longtable}
}

\subsection{Cerrar Ventanas}\label{cerrar-ventanas}

{\def\LTcaptype{none} % do not increment counter
\begin{longtable}[]{@{}ll@{}}
\toprule\noalign{}
Comando & Acción \\
\midrule\noalign{}
\endhead
\bottomrule\noalign{}
\endlastfoot
\texttt{:q} o \texttt{Ctrl+w\ q} & Cerrar ventana actual \\
\texttt{:only} o \texttt{Ctrl+w\ o} & Cerrar todas menos actual \\
\end{longtable}
}

\subsection{Pestañas (Tabs)}\label{pestauxf1as-tabs}

\subsection{Gestión de Pestañas}\label{gestiuxf3n-de-pestauxf1as}

{\def\LTcaptype{none} % do not increment counter
\begin{longtable}[]{@{}ll@{}}
\toprule\noalign{}
Comando & Acción \\
\midrule\noalign{}
\endhead
\bottomrule\noalign{}
\endlastfoot
\texttt{:tabnew} & Nueva pestaña \\
\texttt{:tabe\ archivo} & Abrir archivo en pestaña \\
\texttt{:tabc} & Cerrar pestaña actual \\
\texttt{:tabo} & Cerrar todas las pestañas menos actual \\
\texttt{gt} & Siguiente pestaña \\
\texttt{gT} & Pestaña anterior \\
\texttt{\{n\}gt} & Ir a pestaña n \\
\texttt{:tabs} & Listar pestañas \\
\end{longtable}
}

\section{Buffers}\label{buffers}

\subsection{¿Qué son los Buffers?}\label{quuxe9-son-los-buffers}

Un buffer es un archivo cargado en memoria. Puedes tener múltiples
buffers abiertos.

\subsection{Comandos de Buffer}\label{comandos-de-buffer}

{\def\LTcaptype{none} % do not increment counter
\begin{longtable}[]{@{}ll@{}}
\toprule\noalign{}
Comando & Acción \\
\midrule\noalign{}
\endhead
\bottomrule\noalign{}
\endlastfoot
\texttt{:ls} o \texttt{:buffers} & Listar buffers \\
\texttt{:b\ \{n\}} & Ir a buffer n \\
\texttt{:b\ \{nombre\}} & Ir a buffer por nombre (Tab complete) \\
\texttt{:bn} o \texttt{:bnext} & Siguiente buffer \\
\texttt{:bp} o \texttt{:bprevious} & Buffer anterior \\
\texttt{:bf} & Primer buffer \\
\texttt{:bl} & Último buffer \\
\texttt{:bd} & Eliminar buffer (cerrar archivo) \\
\texttt{:e\ archivo} & Editar archivo (crear buffer) \\
\texttt{:badd\ archivo} & Añadir archivo a buffer list \\
\end{longtable}
}

\subsection{Navegación Rápida}\label{navegaciuxf3n-ruxe1pida}

\begin{Shaded}
\begin{Highlighting}[]
\NormalTok{Ctrl+\^{}      " Alternar entre buffer actual y anterior}
\NormalTok{Ctrl+o      " Saltar a posición anterior}
\NormalTok{Ctrl+i      " Saltar a posición siguiente}
\end{Highlighting}
\end{Shaded}

\section{Búsqueda y Reemplazo}\label{buxfasqueda-y-reemplazo}

\subsection{Comando Substitute}\label{comando-substitute}

\textbf{Sintaxis:}

\begin{Shaded}
\begin{Highlighting}[]
\NormalTok{:[rango]s/patrón/reemplazo/[flags]}
\end{Highlighting}
\end{Shaded}

\subsection{Rangos}\label{rangos}

{\def\LTcaptype{none} % do not increment counter
\begin{longtable}[]{@{}ll@{}}
\toprule\noalign{}
Rango & Descripción \\
\midrule\noalign{}
\endhead
\bottomrule\noalign{}
\endlastfoot
\texttt{\%} & Todo el archivo \\
\texttt{.} & Línea actual \\
\texttt{\$} & Última línea \\
\texttt{\textquotesingle{}\textless{},\textquotesingle{}\textgreater{}}
& Selección visual \\
\texttt{\{n\},\{m\}} & Líneas n a m \\
\texttt{.,+5} & Línea actual + 5 líneas \\
\texttt{g/patrón/} & Líneas que coinciden con patrón \\
\end{longtable}
}

\subsection{Flags}\label{flags}

{\def\LTcaptype{none} % do not increment counter
\begin{longtable}[]{@{}ll@{}}
\toprule\noalign{}
Flag & Descripción \\
\midrule\noalign{}
\endhead
\bottomrule\noalign{}
\endlastfoot
\texttt{g} & Global (todas las ocurrencias en línea) \\
\texttt{c} & Confirmar cada reemplazo \\
\texttt{i} & Case-insensitive \\
\texttt{I} & Case-sensitive \\
\texttt{n} & Reportar número de coincidencias \\
\end{longtable}
}

\subsection{Ejemplos Prácticos}\label{ejemplos-pruxe1cticos}

\begin{Shaded}
\begin{Highlighting}[]
\NormalTok{\# Reemplazar en todo el archivo}
\NormalTok{:\%s/foo/bar/g}

\NormalTok{\# Reemplazar en líneas 10{-}20}
\NormalTok{:10,20s/old/new/g}

\NormalTok{\# Reemplazar con confirmación}
\NormalTok{:\%s/foo/bar/gc}

\NormalTok{\# Reemplazar solo palabras completas}
\NormalTok{:\%s/\textbackslash{}\textless{}foo\textbackslash{}\textgreater{}/bar/g}

\NormalTok{\# Eliminar espacios al final de línea}
\NormalTok{:\%s/\textbackslash{}s\textbackslash{}+$//g}

\NormalTok{\# Añadir ; al final de líneas con texto}
\NormalTok{:\%s/\textbackslash{}(.\textbackslash{}+\textbackslash{})/\textbackslash{}1;/g}

\NormalTok{\# Reemplazar en selección visual}
\NormalTok{:\textquotesingle{}\textless{},\textquotesingle{}\textgreater{}s/old/new/g}

\NormalTok{\# Case{-}insensitive}
\NormalTok{:\%s/error/ERROR/gi}

\NormalTok{\# Eliminar líneas vacías}
\NormalTok{:g/\^{}$/d}

\NormalTok{\# Eliminar líneas que contengan patrón}
\NormalTok{:g/DEBUG/d}

\NormalTok{\# Mantener solo líneas que coincidan}
\NormalTok{:v/KEEP/d}

\NormalTok{\# Duplicar líneas que contengan patrón}
\NormalTok{:g/TODO/t.}

\NormalTok{\# Mover líneas que contengan patrón}
\NormalTok{:g/ERROR/m$}
\end{Highlighting}
\end{Shaded}

\subsection{Caracteres Especiales en
Reemplazo}\label{caracteres-especiales-en-reemplazo}

{\def\LTcaptype{none} % do not increment counter
\begin{longtable}[]{@{}ll@{}}
\toprule\noalign{}
Carácter & Significado \\
\midrule\noalign{}
\endhead
\bottomrule\noalign{}
\endlastfoot
\texttt{\&} & Texto coincidente completo \\
\texttt{\textbackslash{}0} & Texto coincidente completo \\
\texttt{\textbackslash{}1}-\texttt{\textbackslash{}9} & Grupo capturado
1-9 \\
\texttt{\textbackslash{}L} & Siguiente en minúsculas \\
\texttt{\textbackslash{}U} & Siguiente en mayúsculas \\
\texttt{\textbackslash{}l} & Siguiente carácter en minúscula \\
\texttt{\textbackslash{}u} & Siguiente carácter en mayúscula \\
\texttt{\textbackslash{}e} & Fin de \texttt{\textbackslash{}U} o
\texttt{\textbackslash{}L} \\
\end{longtable}
}

\textbf{Ejemplos:}

\begin{Shaded}
\begin{Highlighting}[]
\NormalTok{\# Convertir "lastName" a "last\_name"}
\NormalTok{:\%s/\textbackslash{}([a{-}z]\textbackslash{})\textbackslash{}([A{-}Z]\textbackslash{})/\textbackslash{}1\_\textbackslash{}L\textbackslash{}2/g}

\NormalTok{\# Convertir a mayúsculas primera letra}
\NormalTok{:\%s/\textbackslash{}\textless{}\textbackslash{}(\textbackslash{}w\textbackslash{})/\textbackslash{}u\textbackslash{}1/g}

\NormalTok{\# Invertir orden (swap)}
\NormalTok{:\%s/\textbackslash{}(.*\textbackslash{}), \textbackslash{}(.*\textbackslash{})/\textbackslash{}2 \textbackslash{}1/g}
\end{Highlighting}
\end{Shaded}

\section{Plegado (Folding)}\label{plegado-folding}

\subsection{¿Qué es Folding?}\label{quuxe9-es-folding}

Folding permite ocultar secciones de texto para enfocarse en otras
partes.

\subsection{Métodos de Plegado}\label{muxe9todos-de-plegado}

\begin{Shaded}
\begin{Highlighting}[]
\NormalTok{:set foldmethod=manual    " Manual}
\NormalTok{:set foldmethod=indent    " Por indentación}
\NormalTok{:set foldmethod=syntax    " Por sintaxis}
\NormalTok{:set foldmethod=marker    " Por marcadores}
\NormalTok{:set foldmethod=expr      " Por expresión}
\end{Highlighting}
\end{Shaded}

\subsection{Comandos de Plegado}\label{comandos-de-plegado}

{\def\LTcaptype{none} % do not increment counter
\begin{longtable}[]{@{}ll@{}}
\toprule\noalign{}
Comando & Acción \\
\midrule\noalign{}
\endhead
\bottomrule\noalign{}
\endlastfoot
\texttt{zf\{movimiento\}} & Crear fold (manual) \\
\texttt{zf3j} & Fold 3 líneas abajo \\
\texttt{zfap} & Fold párrafo \\
\texttt{za} & Toggle fold \\
\texttt{zA} & Toggle fold recursivamente \\
\texttt{zo} & Abrir fold \\
\texttt{zO} & Abrir fold recursivamente \\
\texttt{zc} & Cerrar fold \\
\texttt{zC} & Cerrar fold recursivamente \\
\texttt{zd} & Eliminar fold \\
\texttt{zD} & Eliminar folds recursivamente \\
\texttt{zR} & Abrir todos los folds \\
\texttt{zM} & Cerrar todos los folds \\
\texttt{zj} & Mover al siguiente fold \\
\texttt{zk} & Mover al fold anterior \\
\end{longtable}
}

\subsection{Nivel de Plegado}\label{nivel-de-plegado}

\begin{Shaded}
\begin{Highlighting}[]
\NormalTok{:set foldlevel=0     " Cerrar todos}
\NormalTok{:set foldlevel=99    " Abrir todos}
\NormalTok{zm                   " Aumentar foldlevel (cerrar más)}
\NormalTok{zr                   " Disminuir foldlevel (abrir más)}
\end{Highlighting}
\end{Shaded}

\section{Marcas y Saltos}\label{marcas-y-saltos}

\subsection{Marcas (Marks)}\label{marcas-marks}

Guardar posiciones en el archivo.

\subsection{Establecer Marcas}\label{establecer-marcas}

\begin{Shaded}
\begin{Highlighting}[]
\NormalTok{m\{a{-}z\}      " Marca local al buffer}
\NormalTok{m\{A{-}Z\}      " Marca global (entre archivos)}
\end{Highlighting}
\end{Shaded}

\subsection{Ir a Marcas}\label{ir-a-marcas}

{\def\LTcaptype{none} % do not increment counter
\begin{longtable}[]{@{}ll@{}}
\toprule\noalign{}
Comando & Acción \\
\midrule\noalign{}
\endhead
\bottomrule\noalign{}
\endlastfoot
\texttt{\textquotesingle{}\{marca\}} & Ir a línea de marca \\
\texttt{\textasciigrave{}\{marca\}} & Ir a posición exacta \\
\texttt{\textquotesingle{}\textquotesingle{}} & Ir a posición antes del
último salto \\
\texttt{\textasciigrave{}\textasciigrave{}} & Ir a posición exacta antes
del salto \\
\texttt{\textquotesingle{}.} & Ir a última modificación \\
\texttt{\textasciigrave{}.} & Ir a posición exacta de última
modificación \\
\texttt{\textquotesingle{}\^{}} & Ir a posición donde se salió de
Insert \\
\end{longtable}
}

\subsection{Ver Marcas}\label{ver-marcas}

\begin{Shaded}
\begin{Highlighting}[]
\NormalTok{:marks          " Ver todas las marcas}
\NormalTok{:marks abc      " Ver marcas específicas}
\end{Highlighting}
\end{Shaded}

\subsection{Lista de Saltos (Jump
List)}\label{lista-de-saltos-jump-list}

Vim mantiene historial de saltos.

{\def\LTcaptype{none} % do not increment counter
\begin{longtable}[]{@{}ll@{}}
\toprule\noalign{}
Comando & Acción \\
\midrule\noalign{}
\endhead
\bottomrule\noalign{}
\endlastfoot
\texttt{Ctrl+o} & Salto atrás \\
\texttt{Ctrl+i} o \texttt{Tab} & Salto adelante \\
\texttt{:jumps} & Ver lista de saltos \\
\end{longtable}
}

\subsection{Lista de Cambios (Change
List)}\label{lista-de-cambios-change-list}

{\def\LTcaptype{none} % do not increment counter
\begin{longtable}[]{@{}ll@{}}
\toprule\noalign{}
Comando & Acción \\
\midrule\noalign{}
\endhead
\bottomrule\noalign{}
\endlastfoot
\texttt{g;} & Ir a cambio anterior \\
\texttt{g,} & Ir a siguiente cambio \\
\texttt{:changes} & Ver lista de cambios \\
\end{longtable}
}

\section{Comandos Ex}\label{comandos-ex}

\subsection{Comandos de Archivo}\label{comandos-de-archivo}

\begin{Shaded}
\begin{Highlighting}[]
\NormalTok{:w              " Guardar}
\NormalTok{:w archivo      " Guardar como}
\NormalTok{:w!             " Forzar guardar}
\NormalTok{:wa             " Guardar todos}
\NormalTok{:q              " Salir}
\NormalTok{:q!             " Forzar salir}
\NormalTok{:wq o :x        " Guardar y salir}
\NormalTok{:qa             " Salir de todos}
\NormalTok{:qa!            " Forzar salir de todos}
\NormalTok{:e archivo      " Editar archivo}
\NormalTok{:e!             " Recargar archivo}
\NormalTok{:saveas archivo " Guardar como (cambiar nombre)}
\end{Highlighting}
\end{Shaded}

\subsection{Comandos de Shell}\label{comandos-de-shell}

\begin{Shaded}
\begin{Highlighting}[]
\NormalTok{:!comando       " Ejecutar comando shell}
\NormalTok{:r !comando     " Insertar output de comando}
\NormalTok{:w !comando     " Pasar buffer a comando}
\NormalTok{:\%!comando      " Filtrar buffer por comando}
\end{Highlighting}
\end{Shaded}

\textbf{Ejemplos:}

\begin{Shaded}
\begin{Highlighting}[]
\NormalTok{:!ls            " Listar archivos}
\NormalTok{:r !date        " Insertar fecha}
\NormalTok{:\%!sort         " Ordenar líneas}
\NormalTok{:\%!python {-}m json.tool  " Formatear JSON}
\end{Highlighting}
\end{Shaded}

\subsection{Comandos de Ayuda}\label{comandos-de-ayuda}

\begin{Shaded}
\begin{Highlighting}[]
\NormalTok{:help            " Ayuda general}
\NormalTok{:help comando    " Ayuda de comando específico}
\NormalTok{:help \textquotesingle{}opción\textquotesingle{}   " Ayuda de opción}
\NormalTok{:helpgrep patrón " Buscar en ayuda}
\NormalTok{Ctrl+]           " Seguir enlace (en ayuda)}
\NormalTok{Ctrl+t           " Volver (en ayuda)}
\end{Highlighting}
\end{Shaded}

\section{Configuración
(.vimrc/init.vim)}\label{configuraciuxf3n-.vimrcinit.vim}

\subsection{Ubicaciones}\label{ubicaciones}

\textbf{Vim:}

\begin{itemize}
\tightlist
\item
  Linux/Mac: \texttt{\textasciitilde{}/.vimrc}
\item
  Windows: \texttt{\textasciitilde{}/\_vimrc}
\end{itemize}

\textbf{Neovim:}

\begin{itemize}
\tightlist
\item
  Linux/Mac: \texttt{\textasciitilde{}/.config/nvim/init.vim} o
  \texttt{\textasciitilde{}/.config/nvim/init.lua}
\item
  Windows: \texttt{\textasciitilde{}/AppData/Local/nvim/init.vim}
\end{itemize}

\subsection{Configuración Básica}\label{configuraciuxf3n-buxe1sica}

\begin{Shaded}
\begin{Highlighting}[]
\NormalTok{" ============================================}
\NormalTok{" CONFIGURACIÓN BÁSICA}
\NormalTok{" ============================================}

\NormalTok{" Deshabilitar compatibilidad con vi}
\NormalTok{set nocompatible}

\NormalTok{" Habilitar tipo de archivo y plugins}
\NormalTok{filetype plugin indent on}

\NormalTok{" Habilitar sintaxis}
\NormalTok{syntax enable}

\NormalTok{" Codificación}
\NormalTok{set encoding=utf{-}8}
\NormalTok{set fileencoding=utf{-}8}

\NormalTok{" ============================================}
\NormalTok{" INTERFAZ}
\NormalTok{" ============================================}

\NormalTok{" Números de línea}
\NormalTok{set number          " Mostrar números}
\NormalTok{set relativenumber  " Números relativos}

\NormalTok{" Cursor}
\NormalTok{set cursorline      " Resaltar línea actual}
\NormalTok{set scrolloff=8     " Mantener 8 líneas visibles arriba/abajo}

\NormalTok{" Búsqueda}
\NormalTok{set incsearch       " Búsqueda incremental}
\NormalTok{set hlsearch        " Resaltar búsquedas}
\NormalTok{set ignorecase      " Ignorar mayúsculas en búsqueda}
\NormalTok{set smartcase       " Case{-}sensitive si hay mayúsculas}

\NormalTok{" Interfaz}
\NormalTok{set showcmd         " Mostrar comando en progreso}
\NormalTok{set showmode        " Mostrar modo actual}
\NormalTok{set wildmenu        " Autocompletado de comandos mejorado}
\NormalTok{set laststatus=2    " Siempre mostrar statusline}
\NormalTok{set ruler           " Mostrar posición del cursor}

\NormalTok{" Split behavior}
\NormalTok{set splitbelow      " Split horizontal abajo}
\NormalTok{set splitright      " Split vertical a la derecha}

\NormalTok{" ============================================}
\NormalTok{" EDICIÓN}
\NormalTok{" ============================================}

\NormalTok{" Indentación}
\NormalTok{set tabstop=4       " Ancho de tab}
\NormalTok{set shiftwidth=4    " Ancho de indentación}
\NormalTok{set expandtab       " Usar espacios en lugar de tabs}
\NormalTok{set smartindent     " Indentación inteligente}
\NormalTok{set autoindent      " Mantener indentación}

\NormalTok{" Wrap}
\NormalTok{set wrap            " Ajustar líneas largas}
\NormalTok{set linebreak       " Romper en palabras completas}

\NormalTok{" Clipboard}
\NormalTok{set clipboard=unnamedplus  " Usar clipboard del sistema}

\NormalTok{" Backup y swap}
\NormalTok{set nobackup        " No crear backups}
\NormalTok{set noswapfile      " No crear swap files}
\NormalTok{" O guardarlos en directorio específico:}
\NormalTok{" set backupdir=\textasciitilde{}/.vim/backup}
\NormalTok{" set directory=\textasciitilde{}/.vim/swap}

\NormalTok{" Undo persistente}
\NormalTok{set undofile}
\NormalTok{set undodir=\textasciitilde{}/.vim/undo}

\NormalTok{" Mouse}
\NormalTok{set mouse=a         " Habilitar mouse}

\NormalTok{" Tiempo de espera}
\NormalTok{set timeoutlen=500  " Tiempo para secuencias de teclas}

\NormalTok{" ============================================}
\NormalTok{" MAPEOS}
\NormalTok{" ============================================}

\NormalTok{" Líder key}
\NormalTok{let mapleader = " "      " Usar espacio como líder}
\NormalTok{let maplocalleader = "," " Líder local}

\NormalTok{" Guardar con Ctrl+S}
\NormalTok{nnoremap \textless{}C{-}s\textgreater{} :w\textless{}CR\textgreater{}}
\NormalTok{inoremap \textless{}C{-}s\textgreater{} \textless{}Esc\textgreater{}:w\textless{}CR\textgreater{}}

\NormalTok{" Salir de Insert con jj}
\NormalTok{inoremap jj \textless{}Esc\textgreater{}}

\NormalTok{" Navegar entre ventanas más fácil}
\NormalTok{nnoremap \textless{}C{-}h\textgreater{} \textless{}C{-}w\textgreater{}h}
\NormalTok{nnoremap \textless{}C{-}j\textgreater{} \textless{}C{-}w\textgreater{}j}
\NormalTok{nnoremap \textless{}C{-}k\textgreater{} \textless{}C{-}w\textgreater{}k}
\NormalTok{nnoremap \textless{}C{-}l\textgreater{} \textless{}C{-}w\textgreater{}l}

\NormalTok{" Mover líneas}
\NormalTok{nnoremap \textless{}A{-}j\textgreater{} :m .+1\textless{}CR\textgreater{}==}
\NormalTok{nnoremap \textless{}A{-}k\textgreater{} :m .{-}2\textless{}CR\textgreater{}==}
\NormalTok{vnoremap \textless{}A{-}j\textgreater{} :m \textquotesingle{}\textgreater{}+1\textless{}CR\textgreater{}gv=gv}
\NormalTok{vnoremap \textless{}A{-}k\textgreater{} :m \textquotesingle{}\textless{}{-}2\textless{}CR\textgreater{}gv=gv}

\NormalTok{" Mantener selección al indentar}
\NormalTok{vnoremap \textless{} \textless{}gv}
\NormalTok{vnoremap \textgreater{} \textgreater{}gv}

\NormalTok{" Quitar highlight de búsqueda}
\NormalTok{nnoremap \textless{}leader\textgreater{}\textless{}space\textgreater{} :nohlsearch\textless{}CR\textgreater{}}

\NormalTok{" Split rápido}
\NormalTok{nnoremap \textless{}leader\textgreater{}v :vsplit\textless{}CR\textgreater{}}
\NormalTok{nnoremap \textless{}leader\textgreater{}h :split\textless{}CR\textgreater{}}

\NormalTok{" Buffer navigation}
\NormalTok{nnoremap \textless{}leader\textgreater{}bn :bnext\textless{}CR\textgreater{}}
\NormalTok{nnoremap \textless{}leader\textgreater{}bp :bprevious\textless{}CR\textgreater{}}
\NormalTok{nnoremap \textless{}leader\textgreater{}bd :bdelete\textless{}CR\textgreater{}}

\NormalTok{" ============================================}
\NormalTok{" COLORES}
\NormalTok{" ============================================}

\NormalTok{" Tema (requiere instalación)}
\NormalTok{colorscheme desert  " Tema por defecto}

\NormalTok{" True color}
\NormalTok{if has(\textquotesingle{}termguicolors\textquotesingle{})}
\NormalTok{    set termguicolors}
\NormalTok{endif}

\NormalTok{" ============================================}
\NormalTok{" AUTOCOMANDOS}
\NormalTok{" ============================================}

\NormalTok{augroup vimrc}
\NormalTok{    autocmd!}
\NormalTok{    " Volver a última posición al abrir archivo}
\NormalTok{    autocmd BufReadPost * if line("\textquotesingle{}\textbackslash{}"") \textgreater{} 1 \&\& line("\textquotesingle{}\textbackslash{}"") \textless{}= line("$") | exe "normal! g\textquotesingle{}\textbackslash{}"" | endif}
    
\NormalTok{    " Eliminar espacios al final al guardar}
\NormalTok{    autocmd BufWritePre * :\%s/\textbackslash{}s\textbackslash{}+$//e}
    
\NormalTok{    " Resaltar yank}
\NormalTok{    autocmd TextYankPost * silent! lua vim.highlight.on\_yank()}
\NormalTok{augroup END}

\NormalTok{" ============================================}
\NormalTok{" FUNCIONES PERSONALIZADAS}
\NormalTok{" ============================================}

\NormalTok{" Toggle números relativos}
\NormalTok{function! NumberToggle()}
\NormalTok{    if(\&relativenumber == 1)}
\NormalTok{        set norelativenumber}
\NormalTok{    else}
\NormalTok{        set relativenumber}
\NormalTok{    endif}
\NormalTok{endfunction}
\NormalTok{nnoremap \textless{}leader\textgreater{}n :call NumberToggle()\textless{}CR\textgreater{}}
\end{Highlighting}
\end{Shaded}

\subsection{Opciones Útiles}\label{opciones-uxfatiles}

\begin{Shaded}
\begin{Highlighting}[]
\NormalTok{" Persistencia}
\NormalTok{set hidden          " Permitir buffers ocultos con cambios}

\NormalTok{" Rendimiento}
\NormalTok{set lazyredraw      " No redibujar durante macros}
\NormalTok{set updatetime=300  " Tiempo de espera para swap/CursorHold}

\NormalTok{" Signos}
\NormalTok{set signcolumn=yes  " Siempre mostrar columna de signos}

\NormalTok{" Formato}
\NormalTok{set textwidth=80    " Ancho de texto}
\NormalTok{set colorcolumn=80  " Línea visual en columna 80}

\NormalTok{" Completion}
\NormalTok{set completeopt=menuone,noselect,noinsert}

\NormalTok{" Diff}
\NormalTok{set diffopt+=vertical  " Diff vertical}
\end{Highlighting}
\end{Shaded}

\section{Plugins}\label{plugins}

\subsection{Gestores de Plugins}\label{gestores-de-plugins}

\subsection{vim-plug (Recomendado)}\label{vim-plug-recomendado}

\textbf{Instalación:}

\begin{Shaded}
\begin{Highlighting}[]
\CommentTok{\# Vim}
\ExtensionTok{curl} \AttributeTok{{-}fLo}\NormalTok{ \textasciitilde{}/.vim/autoload/plug.vim }\AttributeTok{{-}{-}create{-}dirs} \DataTypeTok{\textbackslash{}}
\NormalTok{    https://raw.githubusercontent.com/junegunn/vim{-}plug/master/plug.vim}

\CommentTok{\# Neovim}
\FunctionTok{sh} \AttributeTok{{-}c} \StringTok{\textquotesingle{}curl {-}fLo "$\{XDG\_DATA\_HOME:{-}$HOME/.local/share\}"/nvim/site/autoload/plug.vim {-}{-}create{-}dirs \textbackslash{}}
\StringTok{       https://raw.githubusercontent.com/junegunn/vim{-}plug/master/plug.vim\textquotesingle{}}
\end{Highlighting}
\end{Shaded}

\textbf{Uso en .vimrc:}

\begin{Shaded}
\begin{Highlighting}[]
\NormalTok{call plug\#begin(\textquotesingle{}\textasciitilde{}/.vim/plugged\textquotesingle{})}

\NormalTok{" Plugins aquí}
\NormalTok{Plug \textquotesingle{}tpope/vim{-}surround\textquotesingle{}}
\NormalTok{Plug \textquotesingle{}junegunn/fzf.vim\textquotesingle{}}
\NormalTok{Plug \textquotesingle{}preservim/nerdtree\textquotesingle{}}

\NormalTok{call plug\#end()}
\end{Highlighting}
\end{Shaded}

\textbf{Comandos:}

\begin{Shaded}
\begin{Highlighting}[]
\NormalTok{:PlugInstall    " Instalar plugins}
\NormalTok{:PlugUpdate     " Actualizar plugins}
\NormalTok{:PlugClean      " Eliminar plugins no usados}
\NormalTok{:PlugStatus     " Ver estado}
\end{Highlighting}
\end{Shaded}

\subsection{Packer (Neovim + Lua)}\label{packer-neovim-lua}

\begin{Shaded}
\begin{Highlighting}[]
\CommentTok{{-}{-} \textasciitilde{}/.config/nvim/lua/plugins.lua}
\ControlFlowTok{return} \FunctionTok{require}\OperatorTok{(}\StringTok{\textquotesingle{}packer\textquotesingle{}}\OperatorTok{).}\NormalTok{startup}\OperatorTok{(}\KeywordTok{function}\OperatorTok{(}\VariableTok{use}\OperatorTok{)}
\NormalTok{    use }\StringTok{\textquotesingle{}wbthomason/packer.nvim\textquotesingle{}}
\NormalTok{    use }\StringTok{\textquotesingle{}tpope/vim{-}surround\textquotesingle{}}
\NormalTok{    use }\StringTok{\textquotesingle{}nvim{-}telescope/telescope.nvim\textquotesingle{}}
\KeywordTok{end}\OperatorTok{)}
\end{Highlighting}
\end{Shaded}

\subsection{Plugins Esenciales}\label{plugins-esenciales}

\subsection{1. Navegación de Archivos}\label{navegaciuxf3n-de-archivos}

\textbf{NERDTree:}

\begin{Shaded}
\begin{Highlighting}[]
\NormalTok{Plug \textquotesingle{}preservim/nerdtree\textquotesingle{}}

\NormalTok{" Mapeo}
\NormalTok{nnoremap \textless{}C{-}n\textgreater{} :NERDTreeToggle\textless{}CR\textgreater{}}
\end{Highlighting}
\end{Shaded}

\textbf{fzf.vim:}

\begin{Shaded}
\begin{Highlighting}[]
\NormalTok{Plug \textquotesingle{}junegunn/fzf\textquotesingle{}, \{ \textquotesingle{}do\textquotesingle{}: \{ {-}\textgreater{} fzf\#install() \} \}}
\NormalTok{Plug \textquotesingle{}junegunn/fzf.vim\textquotesingle{}}

\NormalTok{" Mapeos}
\NormalTok{nnoremap \textless{}C{-}p\textgreater{} :Files\textless{}CR\textgreater{}}
\NormalTok{nnoremap \textless{}leader\textgreater{}b :Buffers\textless{}CR\textgreater{}}
\NormalTok{nnoremap \textless{}leader\textgreater{}g :Rg\textless{}CR\textgreater{}}
\end{Highlighting}
\end{Shaded}

\subsection{2. Edición}\label{ediciuxf3n}

\textbf{vim-surround:}

\begin{Shaded}
\begin{Highlighting}[]
\NormalTok{Plug \textquotesingle{}tpope/vim{-}surround\textquotesingle{}}

\NormalTok{" Comandos:}
\NormalTok{" ysiw" {-} Rodear palabra con "}
\NormalTok{" cs"\textquotesingle{} {-} Cambiar " por \textquotesingle{}}
\NormalTok{" ds" {-} Eliminar "}
\NormalTok{" yss) {-} Rodear línea con ()}
\end{Highlighting}
\end{Shaded}

\textbf{vim-commentary:}

\begin{Shaded}
\begin{Highlighting}[]
\NormalTok{Plug \textquotesingle{}tpope/vim{-}commentary\textquotesingle{}}

\NormalTok{" Comandos:}
\NormalTok{" gcc {-} Comentar/descomentar línea}
\NormalTok{" gc\{movimiento\} {-} Comentar movimiento}
\NormalTok{" gcap {-} Comentar párrafo}
\end{Highlighting}
\end{Shaded}

\textbf{vim-repeat:}

\begin{Shaded}
\begin{Highlighting}[]
\NormalTok{Plug \textquotesingle{}tpope/vim{-}repeat\textquotesingle{}}
\NormalTok{" Permite usar . con plugins}
\end{Highlighting}
\end{Shaded}

\subsection{3. Git}\label{git}

\textbf{vim-fugitive:}

\begin{Shaded}
\begin{Highlighting}[]
\NormalTok{Plug \textquotesingle{}tpope/vim{-}fugitive\textquotesingle{}}

\NormalTok{" Comandos:}
\NormalTok{:Git       " Git status}
\NormalTok{:Git add \%}
\NormalTok{:Git commit}
\NormalTok{:Git push}
\NormalTok{:Gvdiffsplit  " Ver diff}
\end{Highlighting}
\end{Shaded}

\textbf{gitsigns (Neovim):}

\begin{Shaded}
\begin{Highlighting}[]
\NormalTok{Plug \textquotesingle{}lewis6991/gitsigns.nvim\textquotesingle{}}
\end{Highlighting}
\end{Shaded}

\subsection{4. Temas}\label{temas}

\begin{Shaded}
\begin{Highlighting}[]
\NormalTok{Plug \textquotesingle{}gruvbox{-}community/gruvbox\textquotesingle{}}
\NormalTok{Plug \textquotesingle{}morhetz/gruvbox\textquotesingle{}}
\NormalTok{Plug \textquotesingle{}dracula/vim\textquotesingle{}}
\NormalTok{Plug \textquotesingle{}joshdick/onedark.vim\textquotesingle{}}
\NormalTok{Plug \textquotesingle{}arcticicestudio/nord{-}vim\textquotesingle{}}

\NormalTok{colorscheme gruvbox}
\end{Highlighting}
\end{Shaded}

\subsection{5. Statusline}\label{statusline}

\textbf{vim-airline:}

\begin{Shaded}
\begin{Highlighting}[]
\NormalTok{Plug \textquotesingle{}vim{-}airline/vim{-}airline\textquotesingle{}}
\NormalTok{Plug \textquotesingle{}vim{-}airline/vim{-}airline{-}themes\textquotesingle{}}

\NormalTok{let g:airline\_theme=\textquotesingle{}gruvbox\textquotesingle{}}
\NormalTok{let g:airline\_powerline\_fonts = 1}
\end{Highlighting}
\end{Shaded}

\textbf{lualine (Neovim):}

\begin{Shaded}
\begin{Highlighting}[]
\NormalTok{use }\OperatorTok{\{}
    \StringTok{\textquotesingle{}nvim{-}lualine/lualine.nvim\textquotesingle{}}\OperatorTok{,}
    \VariableTok{requires} \OperatorTok{=} \OperatorTok{\{} \StringTok{\textquotesingle{}kyazdani42/nvim{-}web{-}devicons\textquotesingle{}} \OperatorTok{\}}
\OperatorTok{\}}
\end{Highlighting}
\end{Shaded}

\section{Neovim - Características
Especiales}\label{neovim---caracteruxedsticas-especiales}

\subsection{Diferencias con Vim}\label{diferencias-con-vim}

\textbf{Arquitectura:}

\begin{itemize}
\tightlist
\item
  API asíncrona
\item
  Cliente-servidor
\item
  Soporte para UIs remotas
\item
  Jobs y canales nativos
\end{itemize}

\textbf{Características:}

\begin{itemize}
\tightlist
\item
  LSP integrado
\item
  Tree-sitter (parsing AST)
\item
  Lua como lenguaje de primera clase
\item
  Mejor rendimiento
\item
  API moderna
\end{itemize}

\subsection{Configuración en Lua}\label{configuraciuxf3n-en-lua}

\textbf{init.lua:}

\begin{Shaded}
\begin{Highlighting}[]
\CommentTok{{-}{-} \textasciitilde{}/.config/nvim/init.lua}

\CommentTok{{-}{-} Opciones básicas}
\VariableTok{vim}\OperatorTok{.}\VariableTok{opt}\OperatorTok{.}\VariableTok{number} \OperatorTok{=} \KeywordTok{true}
\VariableTok{vim}\OperatorTok{.}\VariableTok{opt}\OperatorTok{.}\VariableTok{relativenumber} \OperatorTok{=} \KeywordTok{true}
\VariableTok{vim}\OperatorTok{.}\VariableTok{opt}\OperatorTok{.}\VariableTok{expandtab} \OperatorTok{=} \KeywordTok{true}
\VariableTok{vim}\OperatorTok{.}\VariableTok{opt}\OperatorTok{.}\VariableTok{tabstop} \OperatorTok{=} \DecValTok{4}
\VariableTok{vim}\OperatorTok{.}\VariableTok{opt}\OperatorTok{.}\VariableTok{shiftwidth} \OperatorTok{=} \DecValTok{4}
\VariableTok{vim}\OperatorTok{.}\VariableTok{opt}\OperatorTok{.}\VariableTok{smartindent} \OperatorTok{=} \KeywordTok{true}
\VariableTok{vim}\OperatorTok{.}\VariableTok{opt}\OperatorTok{.}\VariableTok{wrap} \OperatorTok{=} \KeywordTok{false}
\VariableTok{vim}\OperatorTok{.}\VariableTok{opt}\OperatorTok{.}\VariableTok{swapfile} \OperatorTok{=} \KeywordTok{false}
\VariableTok{vim}\OperatorTok{.}\VariableTok{opt}\OperatorTok{.}\VariableTok{backup} \OperatorTok{=} \KeywordTok{false}
\VariableTok{vim}\OperatorTok{.}\VariableTok{opt}\OperatorTok{.}\VariableTok{undofile} \OperatorTok{=} \KeywordTok{true}
\VariableTok{vim}\OperatorTok{.}\VariableTok{opt}\OperatorTok{.}\VariableTok{hlsearch} \OperatorTok{=} \KeywordTok{false}
\VariableTok{vim}\OperatorTok{.}\VariableTok{opt}\OperatorTok{.}\VariableTok{incsearch} \OperatorTok{=} \KeywordTok{true}
\VariableTok{vim}\OperatorTok{.}\VariableTok{opt}\OperatorTok{.}\VariableTok{termguicolors} \OperatorTok{=} \KeywordTok{true}
\VariableTok{vim}\OperatorTok{.}\VariableTok{opt}\OperatorTok{.}\VariableTok{scrolloff} \OperatorTok{=} \DecValTok{8}
\VariableTok{vim}\OperatorTok{.}\VariableTok{opt}\OperatorTok{.}\VariableTok{signcolumn} \OperatorTok{=} \StringTok{"yes"}
\VariableTok{vim}\OperatorTok{.}\VariableTok{opt}\OperatorTok{.}\VariableTok{updatetime} \OperatorTok{=} \DecValTok{50}
\VariableTok{vim}\OperatorTok{.}\VariableTok{opt}\OperatorTok{.}\VariableTok{colorcolumn} \OperatorTok{=} \StringTok{"80"}

\CommentTok{{-}{-} Leader key}
\VariableTok{vim}\OperatorTok{.}\VariableTok{g}\OperatorTok{.}\VariableTok{mapleader} \OperatorTok{=} \StringTok{" "}
\VariableTok{vim}\OperatorTok{.}\VariableTok{g}\OperatorTok{.}\VariableTok{maplocalleader} \OperatorTok{=} \StringTok{","}

\CommentTok{{-}{-} Mapeos}
\VariableTok{vim}\OperatorTok{.}\VariableTok{keymap}\OperatorTok{.}\NormalTok{set}\OperatorTok{(}\StringTok{"n"}\OperatorTok{,} \StringTok{"\textless{}leader\textgreater{}pv"}\OperatorTok{,} \VariableTok{vim}\OperatorTok{.}\VariableTok{cmd}\OperatorTok{.}\VariableTok{Ex}\OperatorTok{)}
\VariableTok{vim}\OperatorTok{.}\VariableTok{keymap}\OperatorTok{.}\NormalTok{set}\OperatorTok{(}\StringTok{"n"}\OperatorTok{,} \StringTok{"\textless{}C{-}s\textgreater{}"}\OperatorTok{,} \StringTok{":w\textless{}CR\textgreater{}"}\OperatorTok{)}
\VariableTok{vim}\OperatorTok{.}\VariableTok{keymap}\OperatorTok{.}\NormalTok{set}\OperatorTok{(}\StringTok{"i"}\OperatorTok{,} \StringTok{"jj"}\OperatorTok{,} \StringTok{"\textless{}Esc\textgreater{}"}\OperatorTok{)}

\CommentTok{{-}{-} Mover líneas}
\VariableTok{vim}\OperatorTok{.}\VariableTok{keymap}\OperatorTok{.}\NormalTok{set}\OperatorTok{(}\StringTok{"v"}\OperatorTok{,} \StringTok{"J"}\OperatorTok{,} \StringTok{":m \textquotesingle{}\textgreater{}+1\textless{}CR\textgreater{}gv=gv"}\OperatorTok{)}
\VariableTok{vim}\OperatorTok{.}\VariableTok{keymap}\OperatorTok{.}\NormalTok{set}\OperatorTok{(}\StringTok{"v"}\OperatorTok{,} \StringTok{"K"}\OperatorTok{,} \StringTok{":m \textquotesingle{}\textless{}{-}2\textless{}CR\textgreater{}gv=gv"}\OperatorTok{)}

\CommentTok{{-}{-} Mantener cursor centrado}
\VariableTok{vim}\OperatorTok{.}\VariableTok{keymap}\OperatorTok{.}\NormalTok{set}\OperatorTok{(}\StringTok{"n"}\OperatorTok{,} \StringTok{"\textless{}C{-}d\textgreater{}"}\OperatorTok{,} \StringTok{"\textless{}C{-}d\textgreater{}zz"}\OperatorTok{)}
\VariableTok{vim}\OperatorTok{.}\VariableTok{keymap}\OperatorTok{.}\NormalTok{set}\OperatorTok{(}\StringTok{"n"}\OperatorTok{,} \StringTok{"\textless{}C{-}u\textgreater{}"}\OperatorTok{,} \StringTok{"\textless{}C{-}u\textgreater{}zz"}\OperatorTok{)}
\VariableTok{vim}\OperatorTok{.}\VariableTok{keymap}\OperatorTok{.}\NormalTok{set}\OperatorTok{(}\StringTok{"n"}\OperatorTok{,} \StringTok{"n"}\OperatorTok{,} \StringTok{"nzzzv"}\OperatorTok{)}
\VariableTok{vim}\OperatorTok{.}\VariableTok{keymap}\OperatorTok{.}\NormalTok{set}\OperatorTok{(}\StringTok{"n"}\OperatorTok{,} \StringTok{"N"}\OperatorTok{,} \StringTok{"Nzzzv"}\OperatorTok{)}

\CommentTok{{-}{-} Navegar ventanas}
\VariableTok{vim}\OperatorTok{.}\VariableTok{keymap}\OperatorTok{.}\NormalTok{set}\OperatorTok{(}\StringTok{"n"}\OperatorTok{,} \StringTok{"\textless{}C{-}h\textgreater{}"}\OperatorTok{,} \StringTok{"\textless{}C{-}w\textgreater{}h"}\OperatorTok{)}
\VariableTok{vim}\OperatorTok{.}\VariableTok{keymap}\OperatorTok{.}\NormalTok{set}\OperatorTok{(}\StringTok{"n"}\OperatorTok{,} \StringTok{"\textless{}C{-}j\textgreater{}"}\OperatorTok{,} \StringTok{"\textless{}C{-}w\textgreater{}j"}\OperatorTok{)}
\VariableTok{vim}\OperatorTok{.}\VariableTok{keymap}\OperatorTok{.}\NormalTok{set}\OperatorTok{(}\StringTok{"n"}\OperatorTok{,} \StringTok{"\textless{}C{-}k\textgreater{}"}\OperatorTok{,} \StringTok{"\textless{}C{-}w\textgreater{}k"}\OperatorTok{)}
\VariableTok{vim}\OperatorTok{.}\VariableTok{keymap}\OperatorTok{.}\NormalTok{set}\OperatorTok{(}\StringTok{"n"}\OperatorTok{,} \StringTok{"\textless{}C{-}l\textgreater{}"}\OperatorTok{,} \StringTok{"\textless{}C{-}w\textgreater{}l"}\OperatorTok{)}

\CommentTok{{-}{-} Resaltar yank}
\VariableTok{vim}\OperatorTok{.}\VariableTok{api}\OperatorTok{.}\NormalTok{nvim\_create\_autocmd}\OperatorTok{(}\StringTok{"TextYankPost"}\OperatorTok{,} \OperatorTok{\{}
    \VariableTok{group} \OperatorTok{=} \VariableTok{vim}\OperatorTok{.}\VariableTok{api}\OperatorTok{.}\NormalTok{nvim\_create\_augroup}\OperatorTok{(}\StringTok{"highlight\_yank"}\OperatorTok{,} \OperatorTok{\{} \VariableTok{clear} \OperatorTok{=} \KeywordTok{true} \OperatorTok{\}),}
    \VariableTok{callback} \OperatorTok{=} \KeywordTok{function}\OperatorTok{()}
        \VariableTok{vim}\OperatorTok{.}\VariableTok{highlight}\OperatorTok{.}\NormalTok{on\_yank}\OperatorTok{()}
    \KeywordTok{end}\OperatorTok{,}
\OperatorTok{\})}
\end{Highlighting}
\end{Shaded}

\subsection{Tree-sitter}\label{tree-sitter}

Parser incremental para mejor resaltado de sintaxis.

\begin{Shaded}
\begin{Highlighting}[]
\NormalTok{use }\OperatorTok{\{}
    \StringTok{\textquotesingle{}nvim{-}treesitter/nvim{-}treesitter\textquotesingle{}}\OperatorTok{,}
    \VariableTok{run} \OperatorTok{=} \StringTok{\textquotesingle{}:TSUpdate\textquotesingle{}}
\OperatorTok{\}}

\FunctionTok{require}\StringTok{\textquotesingle{}nvim{-}treesitter.configs\textquotesingle{}}\OperatorTok{.}\NormalTok{setup }\OperatorTok{\{}
    \VariableTok{ensure\_installed} \OperatorTok{=} \OperatorTok{\{} \StringTok{"python"}\OperatorTok{,} \StringTok{"javascript"}\OperatorTok{,} \StringTok{"lua"}\OperatorTok{,} \StringTok{"vim"} \OperatorTok{\},}
    \VariableTok{highlight} \OperatorTok{=} \OperatorTok{\{}
        \VariableTok{enable} \OperatorTok{=} \KeywordTok{true}\OperatorTok{,}
    \OperatorTok{\},}
\OperatorTok{\}}
\end{Highlighting}
\end{Shaded}

\subsection{Telescope (Fuzzy Finder)}\label{telescope-fuzzy-finder}

\begin{Shaded}
\begin{Highlighting}[]
\NormalTok{use }\OperatorTok{\{}
    \StringTok{\textquotesingle{}nvim{-}telescope/telescope.nvim\textquotesingle{}}\OperatorTok{,}
    \VariableTok{requires} \OperatorTok{=} \OperatorTok{\{} \StringTok{\textquotesingle{}nvim{-}lua/plenary.nvim\textquotesingle{}} \OperatorTok{\}}
\OperatorTok{\}}

\CommentTok{{-}{-} Mapeos}
\VariableTok{vim}\OperatorTok{.}\VariableTok{keymap}\OperatorTok{.}\NormalTok{set}\OperatorTok{(}\StringTok{\textquotesingle{}n\textquotesingle{}}\OperatorTok{,} \StringTok{\textquotesingle{}\textless{}leader\textgreater{}ff\textquotesingle{}}\OperatorTok{,} \StringTok{\textquotesingle{}:Telescope find\_files\textless{}CR\textgreater{}\textquotesingle{}}\OperatorTok{)}
\VariableTok{vim}\OperatorTok{.}\VariableTok{keymap}\OperatorTok{.}\NormalTok{set}\OperatorTok{(}\StringTok{\textquotesingle{}n\textquotesingle{}}\OperatorTok{,} \StringTok{\textquotesingle{}\textless{}leader\textgreater{}fg\textquotesingle{}}\OperatorTok{,} \StringTok{\textquotesingle{}:Telescope live\_grep\textless{}CR\textgreater{}\textquotesingle{}}\OperatorTok{)}
\VariableTok{vim}\OperatorTok{.}\VariableTok{keymap}\OperatorTok{.}\NormalTok{set}\OperatorTok{(}\StringTok{\textquotesingle{}n\textquotesingle{}}\OperatorTok{,} \StringTok{\textquotesingle{}\textless{}leader\textgreater{}fb\textquotesingle{}}\OperatorTok{,} \StringTok{\textquotesingle{}:Telescope buffers\textless{}CR\textgreater{}\textquotesingle{}}\OperatorTok{)}
\end{Highlighting}
\end{Shaded}

\section{LSP y Autocompletado}\label{lsp-y-autocompletado}

\subsection{LSP en Neovim}\label{lsp-en-neovim}

\textbf{Instalar nvim-lspconfig:}

\begin{Shaded}
\begin{Highlighting}[]
\NormalTok{use }\StringTok{\textquotesingle{}neovim/nvim{-}lspconfig\textquotesingle{}}

\CommentTok{{-}{-} Configurar servidor}
\KeywordTok{local} \VariableTok{lspconfig} \OperatorTok{=} \FunctionTok{require}\OperatorTok{(}\StringTok{\textquotesingle{}lspconfig\textquotesingle{}}\OperatorTok{)}

\CommentTok{{-}{-} Python}
\VariableTok{lspconfig}\OperatorTok{.}\VariableTok{pyright}\OperatorTok{.}\NormalTok{setup}\OperatorTok{\{\}}

\CommentTok{{-}{-} JavaScript/TypeScript}
\VariableTok{lspconfig}\OperatorTok{.}\VariableTok{tsserver}\OperatorTok{.}\NormalTok{setup}\OperatorTok{\{\}}

\CommentTok{{-}{-} Lua}
\VariableTok{lspconfig}\OperatorTok{.}\VariableTok{lua\_ls}\OperatorTok{.}\NormalTok{setup}\OperatorTok{\{}
    \VariableTok{settings} \OperatorTok{=} \OperatorTok{\{}
        \VariableTok{Lua} \OperatorTok{=} \OperatorTok{\{}
            \VariableTok{diagnostics} \OperatorTok{=} \OperatorTok{\{}
                \VariableTok{globals} \OperatorTok{=} \OperatorTok{\{} \StringTok{\textquotesingle{}vim\textquotesingle{}} \OperatorTok{\}}
            \OperatorTok{\}}
        \OperatorTok{\}}
    \OperatorTok{\}}
\OperatorTok{\}}
\end{Highlighting}
\end{Shaded}

\subsection{Mapeos LSP}\label{mapeos-lsp}

\begin{Shaded}
\begin{Highlighting}[]
\CommentTok{{-}{-} Después de attach LSP}
\KeywordTok{local} \VariableTok{opts} \OperatorTok{=} \OperatorTok{\{} \VariableTok{noremap}\OperatorTok{=}\KeywordTok{true}\OperatorTok{,} \VariableTok{silent}\OperatorTok{=}\KeywordTok{true} \OperatorTok{\}}
\VariableTok{vim}\OperatorTok{.}\VariableTok{keymap}\OperatorTok{.}\NormalTok{set}\OperatorTok{(}\StringTok{\textquotesingle{}n\textquotesingle{}}\OperatorTok{,} \StringTok{\textquotesingle{}gD\textquotesingle{}}\OperatorTok{,} \VariableTok{vim}\OperatorTok{.}\VariableTok{lsp}\OperatorTok{.}\VariableTok{buf}\OperatorTok{.}\VariableTok{declaration}\OperatorTok{,} \VariableTok{opts}\OperatorTok{)}
\VariableTok{vim}\OperatorTok{.}\VariableTok{keymap}\OperatorTok{.}\NormalTok{set}\OperatorTok{(}\StringTok{\textquotesingle{}n\textquotesingle{}}\OperatorTok{,} \StringTok{\textquotesingle{}gd\textquotesingle{}}\OperatorTok{,} \VariableTok{vim}\OperatorTok{.}\VariableTok{lsp}\OperatorTok{.}\VariableTok{buf}\OperatorTok{.}\VariableTok{definition}\OperatorTok{,} \VariableTok{opts}\OperatorTok{)}
\VariableTok{vim}\OperatorTok{.}\VariableTok{keymap}\OperatorTok{.}\NormalTok{set}\OperatorTok{(}\StringTok{\textquotesingle{}n\textquotesingle{}}\OperatorTok{,} \StringTok{\textquotesingle{}K\textquotesingle{}}\OperatorTok{,} \VariableTok{vim}\OperatorTok{.}\VariableTok{lsp}\OperatorTok{.}\VariableTok{buf}\OperatorTok{.}\VariableTok{hover}\OperatorTok{,} \VariableTok{opts}\OperatorTok{)}
\VariableTok{vim}\OperatorTok{.}\VariableTok{keymap}\OperatorTok{.}\NormalTok{set}\OperatorTok{(}\StringTok{\textquotesingle{}n\textquotesingle{}}\OperatorTok{,} \StringTok{\textquotesingle{}gi\textquotesingle{}}\OperatorTok{,} \VariableTok{vim}\OperatorTok{.}\VariableTok{lsp}\OperatorTok{.}\VariableTok{buf}\OperatorTok{.}\VariableTok{implementation}\OperatorTok{,} \VariableTok{opts}\OperatorTok{)}
\VariableTok{vim}\OperatorTok{.}\VariableTok{keymap}\OperatorTok{.}\NormalTok{set}\OperatorTok{(}\StringTok{\textquotesingle{}n\textquotesingle{}}\OperatorTok{,} \StringTok{\textquotesingle{}\textless{}C{-}k\textgreater{}\textquotesingle{}}\OperatorTok{,} \VariableTok{vim}\OperatorTok{.}\VariableTok{lsp}\OperatorTok{.}\VariableTok{buf}\OperatorTok{.}\VariableTok{signature\_help}\OperatorTok{,} \VariableTok{opts}\OperatorTok{)}
\VariableTok{vim}\OperatorTok{.}\VariableTok{keymap}\OperatorTok{.}\NormalTok{set}\OperatorTok{(}\StringTok{\textquotesingle{}n\textquotesingle{}}\OperatorTok{,} \StringTok{\textquotesingle{}\textless{}leader\textgreater{}rn\textquotesingle{}}\OperatorTok{,} \VariableTok{vim}\OperatorTok{.}\VariableTok{lsp}\OperatorTok{.}\VariableTok{buf}\OperatorTok{.}\VariableTok{rename}\OperatorTok{,} \VariableTok{opts}\OperatorTok{)}
\VariableTok{vim}\OperatorTok{.}\VariableTok{keymap}\OperatorTok{.}\NormalTok{set}\OperatorTok{(}\StringTok{\textquotesingle{}n\textquotesingle{}}\OperatorTok{,} \StringTok{\textquotesingle{}\textless{}leader\textgreater{}ca\textquotesingle{}}\OperatorTok{,} \VariableTok{vim}\OperatorTok{.}\VariableTok{lsp}\OperatorTok{.}\VariableTok{buf}\OperatorTok{.}\VariableTok{code\_action}\OperatorTok{,} \VariableTok{opts}\OperatorTok{)}
\VariableTok{vim}\OperatorTok{.}\VariableTok{keymap}\OperatorTok{.}\NormalTok{set}\OperatorTok{(}\StringTok{\textquotesingle{}n\textquotesingle{}}\OperatorTok{,} \StringTok{\textquotesingle{}gr\textquotesingle{}}\OperatorTok{,} \VariableTok{vim}\OperatorTok{.}\VariableTok{lsp}\OperatorTok{.}\VariableTok{buf}\OperatorTok{.}\VariableTok{references}\OperatorTok{,} \VariableTok{opts}\OperatorTok{)}
\end{Highlighting}
\end{Shaded}

\subsection{nvim-cmp (Autocompletado)}\label{nvim-cmp-autocompletado}

\begin{Shaded}
\begin{Highlighting}[]
\NormalTok{use }\StringTok{\textquotesingle{}hrsh7th/nvim{-}cmp\textquotesingle{}}
\NormalTok{use }\StringTok{\textquotesingle{}hrsh7th/cmp{-}nvim{-}lsp\textquotesingle{}}
\NormalTok{use }\StringTok{\textquotesingle{}hrsh7th/cmp{-}buffer\textquotesingle{}}
\NormalTok{use }\StringTok{\textquotesingle{}hrsh7th/cmp{-}path\textquotesingle{}}
\NormalTok{use }\StringTok{\textquotesingle{}L3MON4D3/LuaSnip\textquotesingle{}}

\KeywordTok{local} \VariableTok{cmp} \OperatorTok{=} \FunctionTok{require}\StringTok{\textquotesingle{}cmp\textquotesingle{}}

\VariableTok{cmp}\OperatorTok{.}\NormalTok{setup}\OperatorTok{(\{}
    \VariableTok{snippet} \OperatorTok{=} \OperatorTok{\{}
        \VariableTok{expand} \OperatorTok{=} \KeywordTok{function}\OperatorTok{(}\VariableTok{args}\OperatorTok{)}
            \FunctionTok{require}\OperatorTok{(}\StringTok{\textquotesingle{}luasnip\textquotesingle{}}\OperatorTok{).}\NormalTok{lsp\_expand}\OperatorTok{(}\VariableTok{args}\OperatorTok{.}\VariableTok{body}\OperatorTok{)}
        \KeywordTok{end}\OperatorTok{,}
    \OperatorTok{\},}
    \VariableTok{mapping} \OperatorTok{=} \OperatorTok{\{}
        \OperatorTok{[}\StringTok{\textquotesingle{}\textless{}C{-}p\textgreater{}\textquotesingle{}}\OperatorTok{]} \OperatorTok{=} \VariableTok{cmp}\OperatorTok{.}\VariableTok{mapping}\OperatorTok{.}\NormalTok{select\_prev\_item}\OperatorTok{(),}
        \OperatorTok{[}\StringTok{\textquotesingle{}\textless{}C{-}n\textgreater{}\textquotesingle{}}\OperatorTok{]} \OperatorTok{=} \VariableTok{cmp}\OperatorTok{.}\VariableTok{mapping}\OperatorTok{.}\NormalTok{select\_next\_item}\OperatorTok{(),}
        \OperatorTok{[}\StringTok{\textquotesingle{}\textless{}C{-}d\textgreater{}\textquotesingle{}}\OperatorTok{]} \OperatorTok{=} \VariableTok{cmp}\OperatorTok{.}\VariableTok{mapping}\OperatorTok{.}\NormalTok{scroll\_docs}\OperatorTok{({-}}\DecValTok{4}\OperatorTok{),}
        \OperatorTok{[}\StringTok{\textquotesingle{}\textless{}C{-}f\textgreater{}\textquotesingle{}}\OperatorTok{]} \OperatorTok{=} \VariableTok{cmp}\OperatorTok{.}\VariableTok{mapping}\OperatorTok{.}\NormalTok{scroll\_docs}\OperatorTok{(}\DecValTok{4}\OperatorTok{),}
        \OperatorTok{[}\StringTok{\textquotesingle{}\textless{}C{-}Space\textgreater{}\textquotesingle{}}\OperatorTok{]} \OperatorTok{=} \VariableTok{cmp}\OperatorTok{.}\VariableTok{mapping}\OperatorTok{.}\NormalTok{complete}\OperatorTok{(),}
        \OperatorTok{[}\StringTok{\textquotesingle{}\textless{}C{-}e\textgreater{}\textquotesingle{}}\OperatorTok{]} \OperatorTok{=} \VariableTok{cmp}\OperatorTok{.}\VariableTok{mapping}\OperatorTok{.}\NormalTok{close}\OperatorTok{(),}
        \OperatorTok{[}\StringTok{\textquotesingle{}\textless{}CR\textgreater{}\textquotesingle{}}\OperatorTok{]} \OperatorTok{=} \VariableTok{cmp}\OperatorTok{.}\VariableTok{mapping}\OperatorTok{.}\NormalTok{confirm}\OperatorTok{(\{} \FunctionTok{select} \OperatorTok{=} \KeywordTok{true} \OperatorTok{\}),}
    \OperatorTok{\},}
    \VariableTok{sources} \OperatorTok{=} \OperatorTok{\{}
        \OperatorTok{\{} \VariableTok{name} \OperatorTok{=} \StringTok{\textquotesingle{}nvim\_lsp\textquotesingle{}} \OperatorTok{\},}
        \OperatorTok{\{} \VariableTok{name} \OperatorTok{=} \StringTok{\textquotesingle{}buffer\textquotesingle{}} \OperatorTok{\},}
        \OperatorTok{\{} \VariableTok{name} \OperatorTok{=} \StringTok{\textquotesingle{}path\textquotesingle{}} \OperatorTok{\},}
    \OperatorTok{\}}
\OperatorTok{\})}
\end{Highlighting}
\end{Shaded}

\section{Vimscript y Lua}\label{vimscript-y-lua}

\subsection{Vimscript Básico}\label{vimscript-buxe1sico}

\textbf{Variables:}

\begin{Shaded}
\begin{Highlighting}[]
\NormalTok{let variable = "valor"}
\NormalTok{let g:global\_var = 1}
\NormalTok{let l:local\_var = 2}
\NormalTok{let s:script\_var = 3}
\end{Highlighting}
\end{Shaded}

\textbf{Condicionales:}

\begin{Shaded}
\begin{Highlighting}[]
\NormalTok{if condición}
\NormalTok{    " código}
\NormalTok{elseif otra\_condición}
\NormalTok{    " código}
\NormalTok{else}
\NormalTok{    " código}
\NormalTok{endif}
\end{Highlighting}
\end{Shaded}

\textbf{Loops:}

\begin{Shaded}
\begin{Highlighting}[]
\NormalTok{for item in lista}
\NormalTok{    " código}
\NormalTok{endfor}

\NormalTok{while condición}
\NormalTok{    " código}
\NormalTok{endwhile}
\end{Highlighting}
\end{Shaded}

\textbf{Funciones:}

\begin{Shaded}
\begin{Highlighting}[]
\NormalTok{function! MiFuncion(arg1, arg2)}
\NormalTok{    " código}
\NormalTok{    return resultado}
\NormalTok{endfunction}

\NormalTok{" Llamar}
\NormalTok{call MiFuncion(1, 2)}
\end{Highlighting}
\end{Shaded}

\subsection{Lua Básico}\label{lua-buxe1sico}

\textbf{Variables:}

\begin{Shaded}
\begin{Highlighting}[]
\KeywordTok{local} \VariableTok{variable} \OperatorTok{=} \StringTok{"valor"}
\VariableTok{vim}\OperatorTok{.}\VariableTok{g}\OperatorTok{.}\VariableTok{global\_var} \OperatorTok{=} \DecValTok{1}
\end{Highlighting}
\end{Shaded}

\textbf{Condicionales:}

\begin{Shaded}
\begin{Highlighting}[]
\ControlFlowTok{if} \VariableTok{condici}\NormalTok{ó}\VariableTok{n} \ControlFlowTok{then}
    \CommentTok{{-}{-} código}
\ControlFlowTok{elseif} \VariableTok{otra\_condici}\NormalTok{ó}\VariableTok{n} \ControlFlowTok{then}
    \CommentTok{{-}{-} código}
\ControlFlowTok{else}
    \CommentTok{{-}{-} código}
\ControlFlowTok{end}
\end{Highlighting}
\end{Shaded}

\textbf{Loops:}

\begin{Shaded}
\begin{Highlighting}[]
\ControlFlowTok{for} \VariableTok{i} \OperatorTok{=} \DecValTok{1}\OperatorTok{,} \DecValTok{10} \ControlFlowTok{do}
    \CommentTok{{-}{-} código}
\ControlFlowTok{end}

\ControlFlowTok{for} \VariableTok{key}\OperatorTok{,} \VariableTok{value} \KeywordTok{in} \FunctionTok{pairs}\OperatorTok{(}\VariableTok{tabla}\OperatorTok{)} \ControlFlowTok{do}
    \CommentTok{{-}{-} código}
\ControlFlowTok{end}
\end{Highlighting}
\end{Shaded}

\textbf{Funciones:}

\begin{Shaded}
\begin{Highlighting}[]
\KeywordTok{local} \KeywordTok{function}\NormalTok{ mi\_funcion}\OperatorTok{(}\VariableTok{arg1}\OperatorTok{,} \VariableTok{arg2}\OperatorTok{)}
    \CommentTok{{-}{-} código}
    \ControlFlowTok{return} \VariableTok{resultado}
\KeywordTok{end}
\end{Highlighting}
\end{Shaded}

\section{Trucos Avanzados}\label{trucos-avanzados}

\subsection{1. Edición de Múltiples
Archivos}\label{ediciuxf3n-de-muxfaltiples-archivos}

\textbf{Argumentos:}

\begin{Shaded}
\begin{Highlighting}[]
\NormalTok{:args *.js          " Cargar todos los .js}
\NormalTok{:argdo \%s/old/new/g | update  " Reemplazar en todos}
\end{Highlighting}
\end{Shaded}

\textbf{Buffers:}

\begin{Shaded}
\begin{Highlighting}[]
\NormalTok{:bufdo \%s/old/new/g | update}
\end{Highlighting}
\end{Shaded}

\textbf{Quickfix:}

\begin{Shaded}
\begin{Highlighting}[]
\NormalTok{:vimgrep /patrón/ **/*.py   " Buscar en archivos}
\NormalTok{:copen                      " Abrir quickfix list}
\NormalTok{:cn                         " Siguiente}
\NormalTok{:cp                         " Anterior}
\end{Highlighting}
\end{Shaded}

\subsection{2. Edición Avanzada con
Expresiones}\label{ediciuxf3n-avanzada-con-expresiones}

\textbf{Incrementar números:}

\begin{Shaded}
\begin{Highlighting}[]
\NormalTok{Ctrl+a      " Incrementar número bajo cursor}
\NormalTok{Ctrl+x      " Decrementar}

\NormalTok{\# En visual block:}
\NormalTok{g Ctrl+a    " Incrementar secuencialmente}
\end{Highlighting}
\end{Shaded}

\textbf{Evaluar expresiones:}

\begin{Shaded}
\begin{Highlighting}[]
\NormalTok{\# En insert mode:}
\NormalTok{Ctrl+r =2+2\textless{}CR\textgreater{}  " Inserta resultado (4)}
\end{Highlighting}
\end{Shaded}

\subsection{3. Cambiar en Múltiples
Líneas}\label{cambiar-en-muxfaltiples-luxedneas}

\textbf{Visual Block:}

\begin{verbatim}
1. Ctrl+v (visual block)
2. Seleccionar líneas
3. I (insertar al inicio) o A (al final)
4. Escribir
5. Esc (se aplica a todas)
\end{verbatim}

\subsection{4. Formatear Código}\label{formatear-cuxf3digo}

\textbf{Formatear párrafo:}

\begin{Shaded}
\begin{Highlighting}[]
\NormalTok{gqap        " Format paragraph}
\NormalTok{gwap        " Format and keep cursor position}
\end{Highlighting}
\end{Shaded}

\textbf{External formatter:}

\begin{Shaded}
\begin{Highlighting}[]
\NormalTok{:\%!python {-}m json.tool     " JSON}
\NormalTok{:\%!xmllint {-}{-}format {-}      " XML}
\NormalTok{:\% !prettier {-}{-}parser babel  " JavaScript}
\end{Highlighting}
\end{Shaded}

\subsection{5. Diff Files}\label{diff-files}

\begin{Shaded}
\begin{Highlighting}[]
\NormalTok{:vert diffsplit archivo2}
\NormalTok{:diffthis       " Marcar ventana actual para diff}
\NormalTok{:diffoff        " Desactivar diff}
\NormalTok{]c              " Siguiente cambio}
\NormalTok{[c              " Cambio anterior}
\NormalTok{do              " Diff obtain (obtener cambio)}
\NormalTok{dp              " Diff put (poner cambio)}
\end{Highlighting}
\end{Shaded}

\subsection{6. Sesiones}\label{sesiones}

\textbf{Guardar sesión:}

\begin{Shaded}
\begin{Highlighting}[]
\NormalTok{:mksession \textasciitilde{}/sesion.vim}
\end{Highlighting}
\end{Shaded}

\textbf{Cargar sesión:}

\begin{Shaded}
\begin{Highlighting}[]
\NormalTok{:source \textasciitilde{}/sesion.vim}
\NormalTok{\# O al iniciar:}
\NormalTok{vim {-}S \textasciitilde{}/sesion.vim}
\end{Highlighting}
\end{Shaded}

\subsection{7. Spell Checking}\label{spell-checking}

\begin{Shaded}
\begin{Highlighting}[]
\NormalTok{:set spell spelllang=es    " Activar corrector español}
\NormalTok{:set spell spelllang=en    " Inglés}

\NormalTok{]s      " Siguiente error}
\NormalTok{[s      " Error anterior}
\NormalTok{z=      " Sugerencias}
\NormalTok{zg      " Añadir palabra a diccionario}
\NormalTok{zw      " Marcar palabra como mal escrita}
\end{Highlighting}
\end{Shaded}

\subsection{8. Abrir URL}\label{abrir-url}

\begin{Shaded}
\begin{Highlighting}[]
\NormalTok{gx      " Abrir URL bajo cursor en navegador}
\end{Highlighting}
\end{Shaded}

\subsection{9. Ejecutar Línea como
Comando}\label{ejecutar-luxednea-como-comando}

\begin{Shaded}
\begin{Highlighting}[]
\NormalTok{:.!sh   " Ejecutar línea actual como comando shell}
\end{Highlighting}
\end{Shaded}

\subsection{10. Formateo de Tabla}\label{formateo-de-tabla}

\begin{Shaded}
\begin{Highlighting}[]
\NormalTok{\# Seleccionar tabla en visual}
\NormalTok{:!column {-}t     " Alinear columnas}
\end{Highlighting}
\end{Shaded}

\section{Solución de Problemas}\label{soluciuxf3n-de-problemas}

\subsection{Problemas Comunes}\label{problemas-comunes}

\subsection{1. Vim/Neovim Lento}\label{vimneovim-lento}

\textbf{Soluciones:}

\begin{Shaded}
\begin{Highlighting}[]
\NormalTok{" Deshabilitar plugins temporalmente}
\NormalTok{vim {-}{-}noplugin}

\NormalTok{" Ver tiempo de inicio}
\NormalTok{vim {-}{-}startuptime tiempo.log}

\NormalTok{" Lazy loading de plugins}
\NormalTok{Plug \textquotesingle{}plugin/name\textquotesingle{}, \{ \textquotesingle{}on\textquotesingle{}: \textquotesingle{}Comando\textquotesingle{} \}}

\NormalTok{" Reducir updatetime}
\NormalTok{set updatetime=300}
\end{Highlighting}
\end{Shaded}

\subsection{2. Clipboard No Funciona}\label{clipboard-no-funciona}

\textbf{Verificar soporte:}

\begin{Shaded}
\begin{Highlighting}[]
\NormalTok{:echo has(\textquotesingle{}clipboard\textquotesingle{})  " Debe retornar 1}
\end{Highlighting}
\end{Shaded}

\textbf{Solución Ubuntu:}

\begin{Shaded}
\begin{Highlighting}[]
\FunctionTok{sudo}\NormalTok{ apt install vim{-}gtk3}
\CommentTok{\# O para Neovim:}
\FunctionTok{sudo}\NormalTok{ apt install xclip}
\end{Highlighting}
\end{Shaded}

\subsection{3. Colores Incorrectos}\label{colores-incorrectos}

\begin{Shaded}
\begin{Highlighting}[]
\NormalTok{" Habilitar true color}
\NormalTok{set termguicolors}

\NormalTok{" O verificar $TERM}
\NormalTok{echo $TERM  " Debe ser xterm{-}256color o similar}
\end{Highlighting}
\end{Shaded}

\subsection{4. LSP No Funciona (Neovim)}\label{lsp-no-funciona-neovim}

\textbf{Verificar:}

\begin{Shaded}
\begin{Highlighting}[]
\NormalTok{:checkhealth        " Diagnóstico completo}
\NormalTok{:LspInfo            " Info de LSP}

\NormalTok{" Instalar servidor manualmente}
\NormalTok{:LspInstall pyright}
\end{Highlighting}
\end{Shaded}

\subsection{5. Plugins No Se Instalan}\label{plugins-no-se-instalan}

\begin{Shaded}
\begin{Highlighting}[]
\NormalTok{" vim{-}plug}
\NormalTok{:PlugInstall}
\NormalTok{:PlugUpdate}

\NormalTok{" Verificar errores}
\NormalTok{:messages}
\end{Highlighting}
\end{Shaded}

\subsection{Debugging}\label{debugging}

\textbf{Modo verbose:}

\begin{Shaded}
\begin{Highlighting}[]
\NormalTok{:set verbose=9      " Nivel de verbosidad}
\NormalTok{:set verbosefile=\textasciitilde{}/vim\_debug.log}
\end{Highlighting}
\end{Shaded}

\textbf{Ver opciones:}

\begin{Shaded}
\begin{Highlighting}[]
\NormalTok{:set all            " Ver todas las opciones}
\NormalTok{:set opción?        " Ver valor de opción específica}
\end{Highlighting}
\end{Shaded}

\textbf{Ver mapeos:}

\begin{Shaded}
\begin{Highlighting}[]
\NormalTok{:map                " Ver todos los mapeos}
\NormalTok{:nmap               " Mapeos en modo Normal}
\NormalTok{:imap               " Mapeos en modo Insert}
\end{Highlighting}
\end{Shaded}

\section{Referencia Rápida}\label{referencia-ruxe1pida}

\subsection{Atajos Esenciales}\label{atajos-esenciales}

\textbf{Navegación:}

\begin{verbatim}
hjkl        Básico
w/b         Palabra adelante/atrás
0/$         Inicio/fin de línea
gg/G        Inicio/fin de archivo
Ctrl+d/u    Media página
f/F{char}   Find carácter
\end{verbatim}

\textbf{Edición:}

\begin{verbatim}
i/a         Insert/append
o/O         Nueva línea abajo/arriba
x/dd        Eliminar carácter/línea
yy/p        Copiar/pegar
u/Ctrl+r    Undo/redo
.           Repetir
\end{verbatim}

\textbf{Visual:}

\begin{verbatim}
v/V/Ctrl+v  Visual char/line/block
y/d/c       Copiar/eliminar/cambiar
>/<         Indentar/desindentar
\end{verbatim}

\textbf{Búsqueda:}

\begin{verbatim}
/texto      Buscar
n/N         Siguiente/anterior
*/#         Buscar palabra bajo cursor
:%s/old/new/g  Reemplazar
\end{verbatim}

\textbf{Archivos:}

\begin{verbatim}
:w          Guardar
:q          Salir
:e archivo  Abrir
:bn/:bp     Buffer siguiente/anterior
\end{verbatim}

\textbf{Ventanas:}

\begin{verbatim}
:sp/:vsp    Split
Ctrl+w hjkl Navegar
Ctrl+w =    Igualar tamaños
\end{verbatim}

\subsection{Comandos por Modo}\label{comandos-por-modo}

\textbf{Normal Mode:}

\begin{itemize}
\tightlist
\item
  Navegación: \texttt{hjkl}, \texttt{w}, \texttt{b}, \texttt{0},
  \texttt{\$}, \texttt{gg}, \texttt{G}
\item
  Edición: \texttt{i}, \texttt{a}, \texttt{o}, \texttt{x}, \texttt{dd},
  \texttt{yy}, \texttt{p}
\item
  Visual: \texttt{v}, \texttt{V}, \texttt{Ctrl+v}
\item
  Búsqueda: \texttt{/}, \texttt{?}, \texttt{*}, \texttt{\#}, \texttt{n},
  \texttt{N}
\item
  Comandos: \texttt{:}
\end{itemize}

\textbf{Insert Mode:}

\begin{itemize}
\tightlist
\item
  Salir: \texttt{Esc}, \texttt{Ctrl+{[}}
\item
  Autocompletado: \texttt{Ctrl+n/p}, \texttt{Ctrl+x\ Ctrl+o}
\item
  Insertar: \texttt{Ctrl+r\{registro\}}
\end{itemize}

\textbf{Visual Mode:}

\begin{itemize}
\tightlist
\item
  Operadores: \texttt{y}, \texttt{d}, \texttt{c},
  \texttt{\textgreater{}}, \texttt{\textless{}}, \texttt{=}
\item
  Cambiar modo: \texttt{v}, \texttt{V}, \texttt{Ctrl+v}
\item
  Objetos: \texttt{aw}, \texttt{iw}, \texttt{ap}, \texttt{i(},
  \texttt{a"}
\end{itemize}

\textbf{Command Mode:}

\begin{itemize}
\tightlist
\item
  Archivo: \texttt{:w}, \texttt{:q}, \texttt{:e}, \texttt{:wq}
\item
  Búsqueda: \texttt{/}, \texttt{?}
\item
  Rango: \texttt{:\%s}, \texttt{:.},
  \texttt{:\textquotesingle{}\textless{},\textquotesingle{}\textgreater{}}
\end{itemize}

\section{Conclusión}\label{conclusiuxf3n}

\subsection{El Viaje de Aprendizaje}\label{el-viaje-de-aprendizaje}

\textbf{Nivel 1 - Básicos (Semana 1):}

\begin{itemize}
\tightlist
\item
  Modos (Normal, Insert)
\item
  Navegación básica (hjkl)
\item
  Guardar y salir (:wq)
\item
  Insert/append (i, a, o)
\end{itemize}

\textbf{Nivel 2 - Movimientos (Mes 1):}

\begin{itemize}
\tightlist
\item
  Movimientos de palabra (w, b, e)
\item
  Línea (0, \$, \^{})
\item
  Búsqueda (/, ?, *, \#)
\item
  Objetos de texto (iw, aw)
\end{itemize}

\textbf{Nivel 3 - Operadores (Meses 2-3):}

\begin{itemize}
\tightlist
\item
  Eliminar/cambiar/copiar (d, c, y)
\item
  Composición (dw, ciw, yap)
\item
  Visual mode (v, V, Ctrl+v)
\item
  Registros básicos
\end{itemize}

\textbf{Nivel 4 - Avanzado (Meses 4-6):}

\begin{itemize}
\tightlist
\item
  Macros (q, @)
\item
  Múltiples buffers/ventanas
\item
  Búsqueda y reemplazo avanzado
\item
  Configuración personalizada
\end{itemize}

\textbf{Nivel 5 - Maestro (6+ meses):}

\begin{itemize}
\tightlist
\item
  Plugins y ecosistema
\item
  Vimscript/Lua
\item
  LSP y autocompletado
\item
  Workflow completamente personalizado
\end{itemize}

\subsection{Recursos}\label{recursos}

\textbf{Tutoriales:}

\begin{itemize}
\tightlist
\item
  \texttt{vimtutor} (incluido con Vim)
\item
  Vim Adventures (juego)
\item
  Openvim.com
\item
  vim.fandom.com
\end{itemize}

\textbf{Documentación:}

\begin{itemize}
\tightlist
\item
  \texttt{:help} (ayuda integrada)
\item
  vimhelp.org
\item
  neovim.io/doc
\end{itemize}

\textbf{Comunidad:}

\begin{itemize}
\tightlist
\item
  r/vim, r/neovim (Reddit)
\item
  vi.stackexchange.com
\item
  GitHub Discussions
\end{itemize}

\subsection{Tips Finales}\label{tips-finales}

\begin{enumerate}
\def\labelenumi{\arabic{enumi}.}
\tightlist
\item
  \textbf{Practica diariamente}: 15-30 minutos al día
\item
  \textbf{Aprende incrementalmente}: No trates de aprender todo a la vez
\item
  \textbf{Usa vimtutor}: El tutorial interactivo es excelente
\item
  \textbf{Evita plugins al inicio}: Aprende Vim vanilla primero
\item
  \textbf{Configura gradualmente}: Añade opciones según las necesites
\item
  \textbf{Usa el help}: \texttt{:help\ \{tema\}} es tu mejor amigo
\item
  \textbf{Sé paciente}: La curva de aprendizaje vale la pena
\end{enumerate}

\section{Publicaciones Similares}\label{publicaciones-similares}

Si te interesó este artículo, te recomendamos que explores otros blogs y
recursos relacionados que pueden ampliar tus conocimientos. Aquí te dejo
algunas sugerencias:

\begin{enumerate}
\def\labelenumi{\arabic{enumi}.}
\tightlist
\item
  \href{https://chaska-x.netlify.app/operating-system/2017-05-21-comandos-de-informacion-windows/index.pdf}{\faIcon{file-pdf}}
  \href{https://chaska-x.netlify.app/operating-system/2017-05-21-comandos-de-informacion-windows}{Comandos
  De Informacion Windows}
\item
  \href{https://chaska-x.netlify.app/operating-system/2019-06-19-adb/index.pdf}{\faIcon{file-pdf}}
  \href{https://chaska-x.netlify.app/operating-system/2019-06-19-adb}{Adb}
\item
  \href{https://chaska-x.netlify.app/operating-system/2021-08-17-limpieza-y-optimizacion-de-pc/index.pdf}{\faIcon{file-pdf}}
  \href{https://chaska-x.netlify.app/operating-system/2021-08-17-limpieza-y-optimizacion-de-pc}{Limpieza
  Y Optimizacion De Pc}
\item
  \href{https://chaska-x.netlify.app/operating-system/2021-10-21-scrcpy/index.pdf}{\faIcon{file-pdf}}
  \href{https://chaska-x.netlify.app/operating-system/2021-10-21-scrcpy}{Scrcpy}
\item
  \href{https://chaska-x.netlify.app/operating-system/2022-05-12-gestionar-versiones-de-jdk-en-kubuntu/index.pdf}{\faIcon{file-pdf}}
  \href{https://chaska-x.netlify.app/operating-system/2022-05-12-gestionar-versiones-de-jdk-en-kubuntu}{Gestionar
  Versiones De Jdk En Kubuntu}
\item
  \href{https://chaska-x.netlify.app/operating-system/2022-07-21-instalar-tor-browser/index.pdf}{\faIcon{file-pdf}}
  \href{https://chaska-x.netlify.app/operating-system/2022-07-21-instalar-tor-browser}{Instalar
  Tor Browser}
\item
  \href{https://chaska-x.netlify.app/operating-system/2022-08-14-crear-enlaces-duros-o-hard-link-en-linux/index.pdf}{\faIcon{file-pdf}}
  \href{https://chaska-x.netlify.app/operating-system/2022-08-14-crear-enlaces-duros-o-hard-link-en-linux}{Crear
  Enlaces Duros O Hard Link En Linux}
\item
  \href{https://chaska-x.netlify.app/operating-system/2022-09-27-comandos-vim/index.pdf}{\faIcon{file-pdf}}
  \href{https://chaska-x.netlify.app/operating-system/2022-09-27-comandos-vim}{Comandos
  Vim}
\item
  \href{https://chaska-x.netlify.app/operating-system/2023-02-16-guia-de-git-y-github/index.pdf}{\faIcon{file-pdf}}
  \href{https://chaska-x.netlify.app/operating-system/2023-02-16-guia-de-git-y-github}{Guia
  De Git Y Github}
\item
  \href{https://chaska-x.netlify.app/operating-system/2023-05-02-00-primeros-pasos-en-linux/index.pdf}{\faIcon{file-pdf}}
  \href{https://chaska-x.netlify.app/operating-system/2023-05-02-00-primeros-pasos-en-linux}{00
  Primeros Pasos En Linux}
\item
  \href{https://chaska-x.netlify.app/operating-system/2023-06-17-01-introduccion-linux/index.pdf}{\faIcon{file-pdf}}
  \href{https://chaska-x.netlify.app/operating-system/2023-06-17-01-introduccion-linux}{01
  Introduccion Linux}
\item
  \href{https://chaska-x.netlify.app/operating-system/2023-06-18-02-distribuciones-linux/index.pdf}{\faIcon{file-pdf}}
  \href{https://chaska-x.netlify.app/operating-system/2023-06-18-02-distribuciones-linux}{02
  Distribuciones Linux}
\item
  \href{https://chaska-x.netlify.app/operating-system/2023-06-19-03-instalacion-linux/index.pdf}{\faIcon{file-pdf}}
  \href{https://chaska-x.netlify.app/operating-system/2023-06-19-03-instalacion-linux}{03
  Instalacion Linux}
\item
  \href{https://chaska-x.netlify.app/operating-system/2023-06-20-04-administracion-particiones-volumenes/index.pdf}{\faIcon{file-pdf}}
  \href{https://chaska-x.netlify.app/operating-system/2023-06-20-04-administracion-particiones-volumenes}{04
  Administracion Particiones Volumenes}
\item
  \href{https://chaska-x.netlify.app/operating-system/2023-07-01-atajos-de-teclado-y-comandos-para-usar-vim/index.pdf}{\faIcon{file-pdf}}
  \href{https://chaska-x.netlify.app/operating-system/2023-07-01-atajos-de-teclado-y-comandos-para-usar-vim}{Atajos
  De Teclado Y Comandos Para Usar Vim}
\item
  \href{https://chaska-x.netlify.app/operating-system/2024-07-15-instalando-specitify/index.pdf}{\faIcon{file-pdf}}
  \href{https://chaska-x.netlify.app/operating-system/2024-07-15-instalando-specitify}{Instalando
  Specitify}
\item
  \href{https://chaska-x.netlify.app/operating-system/2025-07-10-gestiona-tus-dotfiles-con-gnu-stow/index.pdf}{\faIcon{file-pdf}}
  \href{https://chaska-x.netlify.app/operating-system/2025-07-10-gestiona-tus-dotfiles-con-gnu-stow}{Gestiona
  Tus Dotfiles Con Gnu Stow}
\item
  \href{https://chaska-x.netlify.app/operating-system/2025-12-30-guia-de-ranger/index.pdf}{\faIcon{file-pdf}}
  \href{https://chaska-x.netlify.app/operating-system/2025-12-30-guia-de-ranger}{Guia
  De Ranger}
\item
  \href{https://chaska-x.netlify.app/operating-system/2025-12-31-guia-de-kitty-terminal/index.pdf}{\faIcon{file-pdf}}
  \href{https://chaska-x.netlify.app/operating-system/2025-12-31-guia-de-kitty-terminal}{Guia
  De Kitty Terminal}
\item
  \href{https://chaska-x.netlify.app/operating-system/2025-12-31-guia-de-positron-ide/index.pdf}{\faIcon{file-pdf}}
  \href{https://chaska-x.netlify.app/operating-system/2025-12-31-guia-de-positron-ide}{Guia
  De Positron Ide}
\item
  \href{https://chaska-x.netlify.app/operating-system/2026-04-23-guia-de-rsync/index.pdf}{\faIcon{file-pdf}}
  \href{https://chaska-x.netlify.app/operating-system/2026-04-23-guia-de-rsync}{Guia
  De Rsync}
\end{enumerate}

Esperamos que encuentres estas publicaciones igualmente interesantes y
útiles. ¡Disfruta de la lectura!






\end{document}
